
\documentclass[journal, 10pt, twocolumn]{IEEEtran}

\usepackage[T1]{fontenc}
\usepackage[utf8]{inputenc}
\usepackage{cite}
\usepackage{amsmath,amssymb,amsfonts}
\usepackage{algorithmic}
\usepackage{graphicx}
\usepackage{textcomp}
\usepackage{xcolor}
\usepackage{booktabs}
\usepackage{tabularx}
\usepackage{float}
\usepackage{hyperref}

\begin{document}

\title{Threat-Adaptive Post-Quantum Ratcheting for Blockchain Federated Learning in Healthcare}


\author{
\begin{tabular}{cc}
    \begin{tabular}{c}
        \textbf{Smt. Ch Vijayalakshmi\textsuperscript{1}} \\
        \text{Department of Computer Science and Engineering} \\
        \textit{Chaitanya Bharathi Institute of Technology} \\
        Hyderabad, Telangana, India \\
        vijayalakshmich_cse@cbit.ac.in
    \end{tabular}
    &
    \begin{tabular}{c}
        \textbf{Chevula Mukesh Mun Vardhan\textsuperscript{2}} \\
        \text{Department of Computer Science and Engineering} \\
        \textit{Chaitanya Bharathi Institute of Technology} \\
        Hyderabad, Telangana, India \\
        chevulamukesh@gmail.com
    \end{tabular}\\
    \rule{0pt}{1.5em} \\ 
    \multicolumn{2}{c}{
        \begin{tabular}{c}
        \textbf{Koil Konda Harshith\textsuperscript{3}} \\
        \text{Department of Computer Science and Engineering} \\
        \textit{Chaitanya Bharathi Institute of Technology} \\
        Hyderabad, Telangana, India \\
        harshith.koilkonda@gmail.com
    \end{tabular}
    }
\end{tabular}
}


\maketitle

\begin{abstract}
Post-Quantum Blockchain Federated Learning (PQBFL) is critical towards healthcare AI that is decentralized and secure in the quantum age. One of the most important issues to overcome is the constant nature of the symmetric ratcheting threshold L j, which compels to make a trade-off between post-compromise security and the cost of costly Kyber KEM re-keying operations. This paper presents a \textbf{Threat-Adaptive Ratcheting} mechanism that continuously adjusts $L_j$ in real time based on a composite security signal $t \in [0,1]$ derived from blockchain telemetry via exponential-decay event weighting. Under zero threat, the protocol operates at maximum efficiency ($L_j = L_{max} = 20$); under confirmed attack, it tightens automatically to $L_j = L_{min} = 2$, reducing the post-compromise exposure window by up to $10\times$ and without human intervention. We also present implementation level hardening against timing side-channel attacks and nonce-reuse attacks which are inherited by the underlying PQBFL codebase. This is the first threat-adaptive key rotation scheme that we are aware of as implemented in a post-quantum federated learning ecosystem.
\end{abstract}

\begin{IEEEkeywords}
Federated Learning, Post-Quantum Cryptography, Blockchain, Adaptive Ratcheting, Side-Channel Attacks, Healthcare.
\end{IEEEkeywords}

\section{Introduction}

The advent of Federated Learning (FL) has opened the door to the possibility of heterogeneous devices (smart medical sensors, mobile diagnosis nodes, hospital data centers) cooperatively generating and training massive quantities of Artificial Intelligence (AI) models. This decentralized paradigm is especially essential in areas like health, and finances, where records of participants and clinical information are very sensitive, voluminous, and governed by such schemes as HIPAA and GDPR. High dimensional feature extraction and gradient computation at local scale would guarantee privacy, whereas historical centralized aggregation at the server causes computational, memory, and security constraints. The use of blockchain technology also leads to an additional decentralization of FL, as the central server is substituted with a distributed ledger, which gets rid of the single point of failure. FL (BC-FL) will make the updates of the global model immutable and provide the ability to track reputation transparently by using blockchain, which effectively coordinates the collaborative trust in heterogeneous nodes.
\cite{bdfll23}.

Nevertheless, the current BC-FL systems use solely classical cryptographic primitives (e.g., RSA, Elliptic Curve Cryptography) in order to use digital signatures, establish a session, and communicate safely. With the development of quantum computing, that is, in the context of the threat of Shor algorithm being on a par with polynomies in terms of computational power, the classical schemes will cease to exist altogether. In the absence of quantum-safe accelerators, deep structural cryptosystems will be ineptly broken, not only the privacy of the updated world model will be compromised but also the impartiality of the blockchain itself.
\cite{pqflmdpi}, \cite{beskar}. Therefore, Post-Quantum Federated Learning (PQ-FL) has now become one of the most important research directions, combining NIST-standardized Post-Quantum Cryptography (PQC) algorithms, including ML-KEM (CRYSTALS-Kyber) and ML-DSA (CRYSTALS-Dilithium). By taking the initiative and switching to PQC, companies can adequately counter the attacks of Harvest-Now-Decrypt-Later (HNDL) where attackers today are intercepting classical traffic to decrypt it cryptographically tomorrow.

Under PQ-FL, edge aggregation is important although it has various allied challenges. First, edge devices, which include hospital monitors or IoT nodes, are resource-constrained regarding computation, memory and the availability of energy. The ECDH and ECDSA functionalities of classical elliptic curve cryptography have a tiny energy footprint and very small key sizes (e.g. 32 bytes). By contrast, lattice-based cryptography such as ML-KEM is highly multi-dimensional using multi-dimensional polynomial matrices which need keys in the thousands of bytes size, and computation complexities generally limited by $O(n \log n)$. These complicated tasks can be repeated many times, resulting in energy wastage and thermal throttling.

Secure messaging protocols such as Signal also use ratcheting schemes to sequentially generate encryption keys to offer Forward Secrecy (FS) and Post-Compromise Security (PCS). Continuous asymmetric ratcheting (i.e. swapping heavy PQ public keys) at each epoch of training creates strong propagation latency, quickly exhausting resources deployed in edge-deployed devices. A fixed symmetric ratcheting threshold ($L_j$) can loosely mitigate this by relying on lightweight Hash-based Message Authentication Codes (HMAC) for consecutive rounds. However, this forces a drastic static trade-off: a high $L_j$ weakens PCS (leaving the system completely vulnerable during the long window between KEM resets), while a low $L_j$ exponentially increases overhead, ultimately crippling network throughput \cite{pqbfl1}.

Therefore, there is an urgent missing requirement of a lightweight and scalable side-channel immune, feedback-based, and post-quantum federated system that has dynamic telemetry to trade off overhead limits and real-time cryptographic threats purely through observation \cite{schealthcare}, \cite{sctrustworthyai}. Therefore, there is an urgent missing requirement of a lightweight and scalable side-channel immune, feedback-based, and post-quantum federated system that has dynamic telemetry to trade off overhead limits and real-time cryptographic threats.

\begin{figure}[!t]
\centering
\includegraphics[width=\columnwidth]{fig1.png}
\caption{High-level system architecture of the Adaptive PQBFL framework, illustrating the three-tier interaction between Blockchain Smart Contract, Global Aggregation Server (with embedded \texttt{ThreatMonitor}), and Healthcare Federated Learning Clients.}
\label{fig:system_arch}
\end{figure}

\subsection{Motivation}
Although much has been researched on the topic of data aggregation and fusion in Federated Learning, EFC architectures, and post-quantum analytics, the current approaches do not adequately address the requirements of the post-quantum sensitive systems. Theoretical quantum resistance of conventional PQ-FL schemes with NIST-standardized ML-KEM are based on periodic replacement of keys which are fixed.
 \cite{pqflmdpi}. This static protocol forces an implicit computational bloat; frameworks such as KHPRF-based aggregation and ZKFL-PQ \cite{zkflpq} impose fixed parameter sizes resulting in order of 40-50 percent increases in cost of communication and drastic damage to convergence speed of systems. These systems introduce a huge latency burden to these networks of a quarter or more when forced to operate in dynamic variants of threat. Although the static security models are designed to maintain the deep temporal consistency per se, they fail to proactively offload the unnecessary overheads when the network parameters become normal, resulting in under-optimal throughput limits (less than 5 frames/s with dense processing topology) and excessive thermal throttling of end nodes. \cite{comp1}, \cite{comp3}.

On the other hand, bandwidth-optimal lightweight threshold adjustment schemes have an intrinsic loss of 2030 percent in robust vulnerability defense in high dimensional arrays of features. Adaptive key rotation protocols applied to conventional sensor networks (e.g., Deep Q-Network (DQN)-based frame offloading and RL-based rotation. \cite{rlrotation1}, \cite{qfluav}) On the other hand, bandwidth-optimal lightweight threshold adjustment schemes have an intrinsic loss of 2030 percent in robust vulnerability defense in high dimensional arrays of features. Adaptive key rotation protocols applied to conventional sensor networks (e.g., Deep Q-Network (DQN)-based frame offloading and RL-based rotation. \cite{bdfll23}, \cite{scalebcfl} report spectral is an improvement model that repeatedly interacts with homogeneous, static, scheduling models that collapse significantly to failure (with convergence being up to 30\% slower or accuracy destruction of up to 10\%--12\%), if nodes that are composite report fail to adhere to their commitment, or modify message delay.
 


In addition, the existing security models do not give much attention to hardware execution layer. Implemented systems of FLS in healthcare and remote IoT sensors leave execution footprints which can be mathematically exploited with given hardware side-channel domains \cite{schealthcare}, \cite{sctrustworthyai}. With real, operational networks of nodes, in practice, cryptographic bounds remain static and dynamically spill state variation to induce timing discrepancy arrays with attackers being able to infer 35-40 percent of gradient directions in the world using average observation vectors. These combined efforts regard ratcheting and secure deployment as a one-dimensional problem or an N-dimensional problem of black-box optimization in effect in the absence of a real-time adaptability, a feedback loop of multiple dimensions, and side-channel-reinforced primitives, which is the requirement to have an authentic resilient Post-Quantum FL ecosystem.

\subsection{Our Contribution}
This paper presents a new Threat-Adaptive Ratcheting framework, which clearly responds to the underlying security and performance shortcomings identified with the past EFC and PQ-FL systems. Unlike the statically scheduled architectures that presume key rotation to be a fixed iterative process, we would adaptively add a continuous ThreatMonitor block that observes blockchain telemetries that are underway, viz. failed signature checks, hash mismatch and timing anomaly, and would perform the deep exponential decay modeling to give normalized threat urgency of the network state. This policy helps fight Post-Compromise Security degradation by automatically clamping to maximum security ML-KEM rotations during active being compromised, and automatically unlocking rotations during safe operation. This eliminates directly up to 90 percent of redundant lattice overhead with no mathematical validation bounds impact.

Our architecture is special in its execution hardening of the protocol execution layer, with legacy cryptography with uniform array weighting and physically unsecure evaluations of individual bytes. It aggressively compiles string equality testing to constant-time verification circuits, and could reliably mitigate predictable nonce weaknesses by active CSPRNG matrix arrays. This clearly allows no deterministic key-stream recoveries at the expense of algorithm throughput. The architecture incorporates dynamical threat telemetry, classical Blockchain transparency, Post-Quantum lattice mechanics, and physical side-channel guard, which form the foundation of unexplored base of extremely impervious, scaleable, and adaptable post-quantum federated implementations.

The article addresses the following objectives:
\begin{enumerate}
\item In order to derive a new Threat-Adaptive Ratcheting policy to create a dynamic mapping of power curves eliminating the fixed trade-off between the limits of Post-Compromise Security and the computational overheads of ML-KEM in PQ-FL overlay environments.
\item To implement a \texttt{ThreatMonitor} system to combine local, exponential decay modeling directly converting raw blockchain verification failure and timing errors into a normalized threat composite signal.
\item to orchestrate a conservative dual-threshold ratchet logic, which evaluates effectively exploiting rotation margins mathematically, and operate effectively resourceful in collapsing active vulnerability windows efficiently without centralized mediation demands.
\item To add physical leakages at the implementation-level in a smooth manner across the bottom-level protocol stack, by implementing solid constant-time cryptographic comparative analysis, injecting randomized CSPRNG nonces in making AES-GCM operations, and avoiding zero-pad KDF salt arrays strictly.
\item To actively provide complete experimental validation which takes into account both measures of communication throughput optimization as well as the ability of the active attacker to restrict itself to windows of capability to formally prove mathematical Forward Secrecy constraints which extends classical ECDH as well as ECDSA and lattice ML-KEM parameters.
\end{enumerate}

The remaining parts of the paper are organized as follows. Section II plots the existing body of literature on organized post-quantum federated operations. The necessary theoretical framework of underlying PQBFL sequence mechanisms and our deep Threat-Adaptive cryptographic derivations is defined in Section III. Section IV determines Experimental Results scaling explicit latency standards and mathematical thresholds. Finally, there is the conclusion of the paper, Section V.

\section{Literature Survey}
Post-Quantum Cryptography (PQC) has been converged with Blockchain and Federated Learning (FL) in a wave of more focused academic interest thanks to the imminent commoditization of quantum computation. This section systematically evaluates recent literature spanning post-quantum architectures, blockchain enhancements, Byzantine security, and adaptive key management across IoT and healthcare domains.

\subsection{Post-Quantum Federated Learning}
The growing quantum assault on asymmetric cryptography has incentivized transitioning FL infrastructures towards PQC. Gharavi et al. \cite{gharavi2025} proposed PQBFL, a foundational blockchain-based protocol for FL that integrates post-quantum algorithms to secure multi-participant model sharing against interception and tampering. Expanding to healthcare analytics, Commey et al. \cite{commey2025_pqsbfl} developed PQS-BFL, integrating ML-DSA (Dilithium) to authenticate medical data updates on decentralized smart contracts while surviving quantum adversaries, a concept they further extended in subsequent framework analytics \cite{commey2025_health}. For full-stack privacy, Zhang et al. \cite{zhang2025} outlined private federated deep learning leveraging post-quantum security over edge domains, complementing Gurung et al. \cite{gurung2024} which extensively benchmarked quantum-secure blockchained FL paradigms. Similarly, Li et al. \cite{li2024} demonstrated that shifting local aggregation to post-quantum cryptography materially enhances security without crippling communication overhead. Yang et al. \cite{yang2023_secureagg} advanced an optimized post-quantum secure aggregation architecture aiming to reduce the massive lattice-bound payload sizes typical of Kyber encapsulation. Additional surveys and technical reviews on quantum computing use cases further support this necessary transition \cite{hlatshwayo2026, anon_pdf1, anon_survey_qsafe}.

\subsection{Blockchain-Enhanced Federated Learning}
Blockchain integration offers three key properties for FL trust models: immutability of model metadata, transparent reputation tracking, and decentralization. Sharma et al. \cite{sharma2022} evaluated a distributed, secure, and privacy-preserving computing architecture enabled by blockchain and FL over IoT networks, which aligns with recent security evaluations in the IoE ecosystem \cite{eren2025} and privacy-preserving mechanisms \cite{sameera2024, moore2024}. Nguyen et al. \cite{nguyen2025} formulated a novel blockchain-enabled FL scheme tailored for IoT anomaly detection, ensuring robust malicious node isolation. In medical domains, Wang et al. \cite{wang2026} provided a systematic demonstration of securing FL with blockchain for highly sensitive hospital fields, while Javed et al. \cite{javed2025} applied AI-driven blockchain structures explicitly to electronic health records. Blockchain-enhanced FL also extends significantly to agricultural environments \cite{marzougui2025} and large foundation models \cite{rong2025}, showing versatile cross-domain utilization \cite{bentley2014, anon_crdownload}.

\subsection{Post-Quantum Blockchain and Consensus}
Integrating quantum resistance directly into the underlying blockchain represents a massive hurdle. Yang et al. iteratively detailed the paradigms separating classical from post-quantum and quantum blockchains in comprehensive surveys \cite{yang2024_survey1, yang2024_survey2}. Similarly, Jose et al. \cite{jose2021} outlined consensus algorithms specifically adapted to post-quantum cryptosystems, bridging the gap highlighted by Chen et al. \cite{chen2021} during their construction of quantum-resistant blockchains for smart cities. Further, module lattice-based authentication frameworks effectively authenticate vehicle networks (e.g., autonomous platooning) directly against quantum capabilities \cite{chaudhary2024}, supported by robust migration strategies outlined across international frameworks \cite{elbizri2026}.

\subsection{Security, Privacy, and Adaptive Dynamics}
To counteract malicious intent such as model data poisoning, Kalapaaking et al. \cite{kalapaaking2023} designed a blockchain-based FL framework using Secure Multi-Party Computation (SMPC) to actively verify global models. Saeed et al. \cite{saeed2024} examined broader cybersecurity approaches for IoT systems combining AI and post-quantum cryptography to address inherent architectural heterogeneity. Baseri et al. \cite{baseri2025} explicitly documented the migration strategies and security risk assessment metrics necessary for proactive defense in the evolving quantum era. Adaptive cryptography progressively scales algorithm selection and key sizes. Kappala et al. \cite{kappala2026} introduced a dynamic, quantum-resistant selective encryption approach tailored for resource-limited sensors, validating dynamic threshold modulation. This aligns intimately with practical cryptographic implementations found in Signal-based messaging \cite{dincer2022}, unverified IoT evaluations \cite{lopez2025}, and multi-user secret key generation in wireless channels \cite{aygul2025}, as well as generic distributed architectures \cite{nezhadsistani2025}. 

\subsection{Literature Summary and Comparison}
\label{subsec:lit_comparison}

Table \ref{tab:lit_survey} provides a systematic comparison of all 36 recent systems against our proposed Adaptive PQBFL framework across evaluation dimensions including post-quantum key exchange, blockchain integration, ratcheting strategy, and adaptivity. 

\textbf{Our model is the first to directly combine NIST-standardized Post-Quantum encryption, Blockchain transparency, Threat-Adaptive ratcheting triggers, and physical side-channel defenses within a single deployable FL protocol.}

\begin{table*}[htbp]
\centering
\caption{Comparison of Adaptive PQBFL with Recent PQC–Blockchain–FL Security Frameworks}
\label{tab:lit_survey}

\footnotesize
\renewcommand{\arraystretch}{1.2}

\begin{tabularx}{\textwidth}{l c c c c c X}
\toprule
\textbf{Reference} & \textbf{Year} & \textbf{Post-Quantum} & \textbf{Blockchain} & \textbf{Ratcheting} & \textbf{Adaptive} & \textbf{Primary Contribution} \\
\midrule

Sharma et al. \cite{sharma2022}
& 2022 & No & Yes & None & No
& Blockchain-supported distributed FL for IoT environments \\

Yang et al. \cite{yang2023_secureagg}
& 2023 & Kyber & No & None & No
& Post-quantum secure aggregation protocol for FL \\

Kalapaaking et al. \cite{kalapaaking2023}
& 2023 & No & Yes & None & No
& SMPC-based poisoning attack mitigation in FL \\

Gurung et al. \cite{gurung2024}
& 2024 & Multiple & Yes & None & No
& Performance benchmarking of PQC algorithms in FL \\

Li et al. \cite{li2024}
& 2024 & Lattice-based & No & None & No
& Quantum-resilient secure FL communication model \\

Commey et al. \cite{pqsbfl}
& 2025 & ML-DSA & Yes & None & No
& PQ-secure FL for healthcare analytics \\

Gharavi et al. \cite{gharavi2025}
& 2025 & Kyber & Yes & None & No
& Blockchain-assisted PQ-secure FL coordination \\

Zhang et al. \cite{zhang2025}
& 2025 & Lattice-based & No & None & No
& End-to-end privacy-preserving FL framework \\

Nguyen et al. \cite{nguyen2025}
& 2025 & No & Yes & None & No
& Blockchain-enabled anomaly detection using FL \\

Baseri et al. \cite{baseri2025}
& 2025 & Generic PQC & Yes & None & No
& Risk-aware FL architecture for quantum-era threats \\

Kappala et al. \cite{kappala2026}
& 2026 & Kyber & No & Adaptive & Yes
& Selective adaptive encryption for efficient FL protection \\

Wang et al. \cite{wang2026}
& 2026 & No & Yes & None & No
& Secure medical-data FL infrastructure \\

Chen et al. \cite{chen2021}
& 2021 & Yes & Yes & None & No
& PQ-secure blockchain smart-city architecture \\

Yang et al. \cite{yang2024_survey1}
& 2024 & Yes & Yes & None & No
& Survey of PQ-secure blockchain architectures \\

Chaudhary et al. \cite{chaudhary2024}
& 2024 & Lattice-based & Yes & None & No
& Blockchain authentication using PQ primitives \\

El Bizri et al. \cite{elbizri2026}
& 2026 & Yes & No & None & No
& Migration roadmap toward PQ-secure infrastructures \\

Jose et al. \cite{jose2021}
& 2021 & Yes & Yes & None & No
& Consensus-layer security in blockchain systems \\

Aygul et al. \cite{aygul2025}
& 2025 & No & No & None & No
& Secret-key generation protocol for WSN security \\

Dincer et al. \cite{dincer2022}
& 2022 & No & No & Signal-based & No
& Signal-assisted lightweight cryptographic protection \\

Saeed et al. \cite{saeed2024}
& 2024 & Yes & No & None & No
& AI-integrated PQC protection for IoT ecosystems \\

Marzougui et al. \cite{marzougui2025}
& 2025 & No & Yes & None & No
& Blockchain-enabled agricultural FL framework \\

Javed et al. \cite{javed2025}
& 2025 & No & Yes & None & No
& Secure EHR sharing using blockchain-assisted FL \\

Eren et al. \cite{eren2025}
& 2025 & Yes & Yes & None & No
& Security architecture for Internet-of-Everything using PQC \\

Lopez et al. \cite{lopez2025}
& 2025 & Yes & No & None & No
& Edge-device PQC deployment evaluation \\

Nezhadsistani et al. \cite{nezhadsistani2025}
& 2025 & No & Yes & None & No
& Healthcare FL using blockchain coordination \\

Sameera et al. \cite{sameera2024}
& 2024 & No & Yes & None & No
& Privacy-preserving blockchain-assisted FL \\

Moore et al. \cite{moore2024}
& 2024 & No & Yes & None & No
& Distributed private FL with blockchain integration \\

Hlatshwayo et al. \cite{hlatshwayo2026}
& 2026 & Yes & No & None & No
& Quantum-computing applications in finance security \\

Rong et al. \cite{rong2025}
& 2025 & No & No & None & No
& Federated learning for large-domain model training \\

Commey et al. \cite{commey2025_health}
& 2025 & Yes & Yes & None & No
& Viewer-based PQ-secure healthcare analytics \\

Yang et al. \cite{yang2024_survey2}
& 2024 & Yes & Yes & None & No
& Comparative survey of PQ blockchain systems \\

Bentley et al. \cite{bentley2014}
& 2014 & No & No & None & No
& Collective-behaviour modelling in large-scale systems \\

Anon Survey \cite{anon_survey_qsafe}
& 2014 & Yes & Yes & None & No
& Early quantum-safe infrastructure survey \\

Anon Signature Study \cite{anon_pdf1}
& 2024 & Yes & No & None & No
& Advances in PQ digital signature mechanisms \\

Anon Blockchain Study \cite{anon_crdownload}
& 2024 & No & Yes & None & No
& Exploratory blockchain security framework \\

\midrule

\textbf{Proposed Adaptive PQBFL}
& \textbf{2026}
& \textbf{Kyber}
& \textbf{Yes}
& \textbf{Adaptive}
& \textbf{Yes}
& \textbf{Unified post-quantum secure federated learning framework integrating blockchain coordination with adaptive ratcheting-based selective encryption}

\\
\bottomrule
\end{tabularx}
\end{table*}

\section{Preliminaries}

\begin{figure*}
\centering
% Generate from Mermaid: mermaid_diagrams.md #2
\includegraphics[width=2\columnwidth]{fig2 (1).png}
\caption{State transition diagram of the four-phase PQBFL protocol lifecycle. Phase~3 iterates symmetrically until the adaptive ratchet policy triggers a Phase~4 asymmetric key rotation, establishing a new session epoch $j+1$.}
\label{fig:protocol_lifecycle}
\end{figure*}

This section is a presentation of the background definitions, architectural, and symbolic notations used in the proposed Threat-Adaptive Post-Quantum Blockchain Federated Learning (PQBFL) system of secure healthcare analytics.

Let $\mathcal{N} = \{N_1, N_2, \dots, N_M\}$ denote the set of $M$ distributed healthcare edge devices (participants) collaboratively training a global machine learning model, anchored to a decentralized Smart Contract ledger $\mathcal{L}$. Each edge device $N_i$ maintains a localized, highly sensitive, and strictly private dataset denoted by $\mathcal{D}_i = \{(x_{i,1}, y_{i,1}), \dots, (x_{i,k}, y_{i,k})\}$ containing $k$ diagnostic samples. The objective is to compute a global parameter matrix $\mathcal{W}$ minimizing the empirical loss function $F(\mathcal{W}) = \sum_{i=1}^M p_i f_i(\mathcal{W}, \mathcal{D}_i)$ where $p_i = |\mathcal{D}_i| / \sum |\mathcal{D}_j|$ acts as the relative data weighting.

Each $N_i$ training round $r$ requires local gradient extraction $\Delta_{i,r}$ using stochastic gradient descent (SGD). Rather than transmitting raw parameters, all communication is symmetrically constrained through Authenticated Encryption with Associated Data (AEAD) $C_{i,r} = \text{Enc}(K_{i,j}, \Delta_{i,r})$, where $K_{i,j} \in \mathcal{K}_S$ represents the volatile symmetric model key dynamically derived for round $i$ inside session window $j$.

The Post-Quantum cryptographic constraints introduce multi-dimensional key spaces. The aggregator server $S$ and participants $\mathcal{N}$ formulate structural key generations spanning dual-domain architectures. Specifically, the classical domain constructs Elliptic Curve Diffie-Hellman matrices $(epk, esk) \in \mathcal{C}$ operating natively on the \textit{secp256k1} boundaries, while the quantum-resistant domain executes ML-KEM (Kyber) polynomial sampling yielding massive key vectors $(kpk, ksk) \in R_q^{k \times k}$. The unified session root key is mapped via a fusion operation $\Phi(\cdot)$ such that the resulting root state is $RK_j = \Phi(S_{kem}, S_{ecdh})$, completely shielding the topology from distinct algorithm compromise vectors.

The entire key ratcheting lifecycle is governed by an adaptive evaluation process modeled over continuous blockchain metadata. The \texttt{ThreatMonitor} module isolates system states evaluating localized security verification boundaries. Let $\mathcal{E}_t = \{e_1, e_2, \dots, e_N\}$ denote the continuous stream of operational disruption anomalies (e.g., hash collisions, signature verification failures, and latency spikes) observed strictly at time $t$. The localized state telemetry uniquely scales based on active severity limits $s_i$, converting event sets into a normalized real-time evaluation vector $\Theta(t) \in [0, 1]$.

Based strictly on this live evaluation, the active policy function $\pi(\Theta)$ inherently dictates the maximum symmetric derivation limit $L_j$ for the current cryptographic session, mapping the perceived urgency onto rotation dimensions $L_j = \pi(\Theta)$. Unlike traditional cryptographic protocols operating strictly on rigid, manually-scaled timeframes ($L_{fixed}$), the proposed architecture evaluates algorithmic shifts whenever $\Theta(t)$ exceeds specific baseline limits. Whenever localized security boundaries fall below tolerance bounds or threat evaluations spike above limit $\mathcal{T}_{attack}$, the dynamic weight scaling actively overrides the global state, instantly shrinking the post-compromise attack window $L_j \to L_{min}$. This effectively forces an immediate asymmetric ratification request $AsymReq_{S \to \mathcal{N}}$, terminating exposed keys instantly while allowing the framework to safely maximize throughput $L_j \to L_{max}$ entirely during benign baseline analytics stages.


\subsection{Notation and Terminology}
Table~\ref{tab:notation} summarises the most important symbols and abbreviated terms used throughout this paper.

\begin{table}[htbp]
\centering
\caption{Important Notation and Terminology}
\label{tab:notation}
\begin{tabularx}{\columnwidth}{l X}
\toprule
\textbf{Symbol} & \textbf{Definition} \\
\midrule
$N_i, S, \mathcal{L}$ & Edge device $i$, Aggregation server, Blockchain ledger \\
$\mathcal{W}, \Delta_{i,r}$ & Global model matrix, Local model gradient at round $r$ \\
$(kpk, ksk)$ & ML-KEM public / secret key pair \\
$(epk, esk)$ & ECDH ephemeral public / secret key pair \\
$(pk, sk)$ & ECDSA signing public / secret key pair \\
$RK_j$ & Hybrid root key for session epoch $j$ \\
$CK_{i,j}$ & Symmetric chain key at round $i$ in epoch $j$ \\
$K_{i,j}$ & AES-GCM model-encryption key \\
$L_j$ & Adaptive symmetric ratchet threshold for epoch $j$ \\
$\mathcal{E}_t$ & Blockchain security anomaly events at time $t$ \\
$\Theta(t)$ & Normalized composite threat level $\in [0,1]$ \\
$\pi(\Theta)$ & Adaptive policy function mapping $\Theta$ to $L_j$ \\
$AsymReq$ & Server-initiated asymmetric re-key request \\
$Tx_r, Tx_p, Tx_u$ & Blockchain transactions (Reg., Publish, Update) \\
HNDL & Harvest-Now-Decrypt-Later attack \\
PCS & Post-Compromise Security \\
\bottomrule
\end{tabularx}
\end{table}


\section{Methodology}
In this section, the hierarchical cryptographic architecture of Threat-Adaptive Post-Quantum Blockchain Federated Learning (PQBFL) is provided throughout the distributed healthcare topology. The general trend of the development of our adaptive system based on DRL-inspired frameworks as opposed to conventional fixed windows of key exposure demonstrates a paradigm shift between inflexible key exposure windows and decentralized feedback-based coordination. Although the baseline method is based on predetermined symmetric thresholds as well as centralized scheduling, our offered system presents distributed telemetry intelligence on each node, real-time threat adaptation, and dual-domain mathematical coupling that contribute to increased responsiveness, scaling, and post-compromise security in heterogeneous IoMulT healthcare context.

The flow of architecture is aimed at the optimal use of secure aggregation of real-time gradient streams of models with anchored feedback in blockchains. The framework involves four major operational stages, including on-chain registration, hybrid session establishment, iterative secure training with symmetric ratcheting and threat-adaptive asymmetric ratchet triggering, synchronized by a ThreatMonitor module and a two-way cryptographic feedback scheme. This architecture presents adaptive coordination at edge and server layers, with ECDSA-based authentication and Kyber-based encapsulation at the cryptographic layer to allow key rotation selection at the cryptographic layer, and ThreatMonitor-based evaluation and dual-threshold logic at the policy layer to provide multiscale security adaptation under heterogeneous and real-time adversarial conditions.

\subsection{Foundational Cryptographic Primitives}
PQBFL framework is based on three fundamental cryptographic algorithms, the mathematical soundness of which as well as the bounds of their computations we establish mathematically, below.


\subsubsection{Elliptic Curve Digital Signature Algorithm (ECDSA)}
To authenticate participants against the blockchain smart contract without exposing private keys, we employ ECDSA over the \textit{secp256k1} curve. All parties agree on domain parameters $(p, a, b, G, n, h)$, where $p$ defines the finite field $\mathbb{F}_p$, the curve equation satisfies $y^2 \equiv x^3 + ax + b \pmod{p}$ with $4a^3 + 27b^2 \not\equiv 0 \pmod{p}$, and $G = (x_G, y_G)$ is the generator point of prime order $n$.

\textbf{Key Generation.} Each participant selects a random private key $d \in [1, n-1]$ and computes the public key via scalar multiplication:
\begin{equation}
Q = d \cdot G = (x_Q, y_Q)
\label{eq:ecdsa_keygen}
\end{equation}
The key pair is defined by Eq.~\eqref{eq:ecdsa_keygen}. The security rests on the Elliptic Curve Discrete Logarithm Problem (ECDLP): given $Q$ and $G$, recovering $d$ is computationally infeasible.

\textbf{Signature Generation.} To sign a message $m$, the signer computes:
\begin{equation}
e = \text{SHA-256}(m), \quad z = \text{leftmost } L_n \text{ bits of } e
\label{eq:ecdsa_hash}
\end{equation}
The hash truncation in Eq.~\eqref{eq:ecdsa_hash} extracts the $L_n$ most-significant bits of $e$. A per-signature nonce $k \xleftarrow{\$} [1, n-1]$ is sampled uniformly at random. Using scalar multiplication on the curve, the commitment point is computed per Eq.~\eqref{eq:ecdsa_r}:
\begin{equation}
P = k \cdot G = (x_1, y_1), \quad r = x_1 \bmod n
\label{eq:ecdsa_r}
\end{equation}
If $r = 0$, a new $k$ is selected. The second signature component (Eq.~\eqref{eq:ecdsa_s}) requires the modular multiplicative inverse $k^{-1} \bmod n$ via the Extended Euclidean Algorithm:
\begin{equation}
s = k^{-1}(z + r \cdot d) \bmod n
\label{eq:ecdsa_s}
\end{equation}
The output signature is the tuple $(r, s)$.

\textbf{Signature Verification.} Given message $m$, signature $(r, s)$, and public key $Q$, the verifier checks $r, s \in [1, n-1]$, then computes the scalar intermediates per Eq.~\eqref{eq:ecdsa_verify_scalars} and the verification point per Eq.~\eqref{eq:ecdsa_verify_point}:
\begin{equation}
\begin{split}
w &= s^{-1} \bmod n, \quad u_1 = z \cdot w \bmod n, \\
u_2 &= r \cdot w \bmod n
\end{split}
\label{eq:ecdsa_verify_scalars}
\end{equation}
\begin{equation}
P_{ver} = u_1 \cdot G + u_2 \cdot Q = (x_p, y_p)
\label{eq:ecdsa_verify_point}
\end{equation}
The signature is valid if and only if $r' = x_p \bmod n = r$.

\textbf{Theorem 1 (ECDSA Correctness Proof).} If the signature $(r,s)$ was generated correctly, verification always succeeds.

\textit{Proof.} Begin with the signing equation:
\begin{equation}
\begin{split}
s &\equiv k^{-1}(z + rd) \pmod{n} \\
\implies k &\equiv s^{-1}(z + rd) \pmod{n}
\end{split}
\label{eq:ecdsa_proof_step1}
\end{equation}
The verification point expands as:
\begin{equation}
\begin{split}
P_{ver} &= u_1G + u_2Q \\
        &= (zs^{-1})G + (rs^{-1})(dG)
\end{split}
\label{eq:ecdsa_proof_step2}
\end{equation}
Factoring out $s^{-1}$ and $G$ by the distributive property of scalar multiplication yields Eq.~\eqref{eq:ecdsa_proof_step3}:
\begin{equation}
P_{ver} = s^{-1}(z + rd) \cdot G
\label{eq:ecdsa_proof_step3}
\end{equation}
Since $k \equiv s^{-1}(z + rd) \pmod{n}$ from Eq.~\eqref{eq:ecdsa_proof_step1}, substituting into Eq.~\eqref{eq:ecdsa_proof_step3} yields $P_{ver} = kG$, which is the exact same point computed during signing via Eq.~\eqref{eq:ecdsa_r}. Therefore $x_p = x_1$, and $r' = r$. $\hfill$

\textbf{Time Complexity Analysis.} The computational bottleneck is scalar multiplication ($k \times G$) as used in Eq.~\eqref{eq:ecdsa_r}. Using the Double-and-Add algorithm across $b$ bits of the scalar $k$, this requires approximately $b$ point doublings and $b/2$ point additions. Each point operation costs $\mathcal{O}(M(b))$ field multiplications, where $M(b) = \mathcal{O}(b^{1.58})$ under Karatsuba. The signature generation and verification complexities are shown in Eq.~\eqref{eq:ecdsa_sign_complexity} and Eq.~\eqref{eq:ecdsa_verify_complexity}, respectively:
\begin{equation}
\begin{split}
T_{\text{ECDSA-Sign}} &= T_{SM} + T_{Inv} \\
                      &= \mathcal{O}(b^3) + \mathcal{O}(b^2) \approx \mathcal{O}(b^3)
\end{split}
\label{eq:ecdsa_sign_complexity}
\end{equation}
\begin{equation}
\begin{split}
T_{\text{ECDSA-Verify}} &= 2 \cdot T_{SM} + T_{Inv} \\
                        &= 2\mathcal{O}(b^3) + \mathcal{O}(b^2) \approx \mathcal{O}(b^3)
\end{split}
\label{eq:ecdsa_verify_complexity}
\end{equation}
For fixed bit-length $b = 256$ (secp256k1), all operations execute in constant time $\mathcal{O}(1)$.

\begin{figure}[!t]
\hrulefill
\vspace{-0.2cm}
\begin{center}
\textbf{Algorithm 3:} ECDSA Signature Generation and Verification
\end{center}
\vspace{-0.2cm}
\hrulefill
\begin{enumerate}
    \item \textbf{Input:} Message $m$, Private key $d$, Public key $Q$, Curve parameters $(p, a, b, G, n, h)$
    \item \textbf{Output:} Signature $(r, s)$ or \textsc{Valid}/\textsc{Invalid}
    \item[]
    \item[] \textbf{--- Signature Generation ---}
    \item $e \leftarrow \text{SHA-256}(m)$
    \item $z \leftarrow \text{leftmost } L_n \text{ bits of } e$
    \item \textbf{repeat}
    \item \quad $k \xleftarrow{\$} [1, n-1]$ \quad \textit{// CSPRNG nonce, must never repeat}
    \item \quad $P \leftarrow k \cdot G$ \quad \textit{// Double-and-Add scalar mult.}
    \item \quad $r \leftarrow x_P \bmod n$
    \item \quad $k^{-1} \leftarrow \text{ExtEuclidean}(k, n)$
    \item \quad $s \leftarrow k^{-1}(z + r \cdot d) \bmod n$
    \item \textbf{until} $r \neq 0$ \textbf{and} $s \neq 0$
    \item \textbf{return} $(r, s)$
    \item[]
    \item[] \textbf{--- Signature Verification ---}
    \item \textbf{if} $r \notin [1, n-1]$ \textbf{or} $s \notin [1, n-1]$ \textbf{then return} \textsc{Invalid}
    \item $e \leftarrow \text{SHA-256}(m)$; \quad $z \leftarrow \text{leftmost } L_n \text{ bits of } e$
    \item $w \leftarrow s^{-1} \bmod n$
    \item $u_1 \leftarrow z \cdot w \bmod n$; \quad $u_2 \leftarrow r \cdot w \bmod n$
    \item $P_{ver} \leftarrow u_1 \cdot G + u_2 \cdot Q$ \quad \textit{// Shamir dual mult.}
    \item \textbf{if} $P_{ver} = \mathcal{O}$ \textbf{then return} \textsc{Invalid}
    \item $r' \leftarrow x_{P_{ver}} \bmod n$
    \item \textbf{return} \textsc{Valid} \textbf{if} $r' = r$, \textsc{Invalid} otherwise
\end{enumerate}
\hrulefill
\end{figure}

\subsubsection{Elliptic Curve Diffie-Hellman (ECDH)}
In the case of the classical layer in our hybrid key agreement, ECDH sets up a common secret between the local side (Bob) and a subject (Alice) using a channel (insecure) that is not secure. The domain parameters of both parties are identical and are given by $(p, a, b, G, n, h)$.

\textbf{Key Exchange Protocol.} The server generates a random private key $esk_S \in [1, n-1]$ and computes $epk_S = esk_S \cdot G$. The participant generates $esk_N \in [1, n-1]$ and computes $epk_N = esk_N \cdot G$. Public keys are exchanged over authenticated channels. Each party independently computes the shared curve point:
\begin{equation}
\begin{split}
S_{\text{server}} &= esk_S \times epk_N, \\
S_{\text{node}} &= esk_N \times epk_S
\end{split}
\label{eq:ecdh_shared}
\end{equation}

\textbf{Theorem 2 (ECDH Shared Secret Alignment).} Both parties compute the identical shared point.

\textit{Proof.} Substitute the public key definitions:
\begin{equation}
\begin{split}
S_{\text{server}} &= esk_S \times (esk_N \times G) \\
                  &= (esk_S \cdot esk_N) \times G
\end{split}
\label{eq:ecdh_proof1}
\end{equation}
\begin{equation}
\begin{split}
S_{\text{node}} &= esk_N \times (esk_S \times G) \\
                &= (esk_N \cdot esk_S) \times G
\end{split}
\label{eq:ecdh_proof2}
\end{equation}
By the associativity and commutativity of scalar multiplication in the cyclic group generated by $G$, Eq.~\eqref{eq:ecdh_proof1} and Eq.~\eqref{eq:ecdh_proof2} both reduce to $(esk_S \cdot esk_N) \times G = (esk_N \cdot esk_S) \times G$. Therefore $S_{\text{server}} = S_{\text{node}}$. The ECDH shared secret is defined per Eq.~\eqref{eq:ecdh_kdf} as the x-coordinate $x_S$ of the agreed point, processed through a KDF:
\begin{equation}
S_{ecdh} = \text{KDF}(x_S)
\label{eq:ecdh_kdf}
\end{equation}
This completes the proof. $\hfill$

\textbf{Time Complexity.} Each party performs exactly two scalar multiplications. The total time complexity per party is $T_{\text{ECDH}} = 2 \cdot T_{SM} = 2 \cdot \mathcal{O}(b^3) = \mathcal{O}(b^3)$, constant for fixed $b = 256$.

\subsubsection{CRYSTALS-Kyber Key Encapsulation Mechanism (ML-KEM)}
The post-quantum security layer utilizes CRYSTALS-Kyber, whose hardness is founded on the Module Learning With Errors (M-LWE) problem. All operations occur in the polynomial ring $R_q = \mathbb{Z}_q[X]/(X^{256}+1)$, where $q = 3329$ for Kyber-768 and any occurrence of $X^{256}$ is replaced by $-1$.

\textbf{The M-LWE Problem.} Let $A \in R_q^{k \times k}$ be a public matrix ($k = 3$ for Kyber-768). A secret vector $s \in R_q^k$ and error vector $e \in R_q^k$ are sampled from a Centered Binomial Distribution with small coefficients $\in \{-1, 0, 1\}$. The resulting M-LWE instance (Eq.~\eqref{eq:mlwe}) forms the public key vector:
\begin{equation}
t = As + e
\label{eq:mlwe}
\end{equation}
Given $(A, t)$, recovering $s$ is computationally infeasible. This forms the bedrock of Kyber's security.

\textbf{Key Generation.} The server generates a 32-byte seed $\rho$ and expands it via the eXtendable-Output Function SHAKE-128 to deterministically construct $A$:
\begin{equation}
\begin{split}
A_{i,j} &= \text{Parse}(\text{SHAKE-128}(\rho \,\|\, i \,\|\, j)), \\
& \quad \forall\, i,j \in \{0, \ldots, k-1\}
\end{split}
\label{eq:kyber_matrix}
\end{equation}
The matrix $A$ is generated via Eq.~\eqref{eq:kyber_matrix}; small secret and error vectors $(s, e)$ are sampled from a CBD. The public key is $pk_{kem} = (A, t)$ and secret key is $sk_{kem} = s$.

\textbf{Encapsulation.} Given the public key $pk_{kem}$, the encapsulating party computes ciphertext components (Eq.~\eqref{eq:kyber_encap1}):
\begin{equation}
u = A^T r + e_1, \quad v' = t^T r + e_2
\label{eq:kyber_encap1}
\end{equation}
where $r, e_1, e_2$ are derived from a random 256-bit message $m$. The message is encoded into a polynomial via $m_{poly} = m \cdot \lfloor q/2 \rfloor$, and the full ciphertext with shared secret is defined in Eq.~\eqref{eq:kyber_encap2}:
\begin{equation}
\begin{split}
v &= v' + m_{poly}, \quad ct = (u, v), \\
SS_k &= \text{KDF}(m, H(ct))
\end{split}
\label{eq:kyber_encap2}
\end{equation}

\textbf{Decapsulation.} The decapsulating party computes $v' = s^T u$ and recovers the noisy message polynomial per Eq.~\eqref{eq:kyber_decap1}:
\begin{equation}
m'_{poly} = v - s^T u
\label{eq:kyber_decap1}
\end{equation}

\textbf{Theorem 3 (Kyber Decapsulation Correctness).} The decapsulated message polynomial equals the original message plus negligible noise.

\textit{Proof.} Expanding $m'_{poly}$:
\begin{equation}
\begin{aligned}
m'_{poly} &= v - s^T u \\
          &= (v' + m_{poly}) - s^T(A^T r + e_1) \\
          &= (t^T r + e_2 + m_{poly}) - (s^T A^T r + s^T e_1) \\
          &= m_{poly} + \underbrace{(e^T r + e_2 - s^T e_1)}_{\text{small noise } \nu}
\end{aligned}
\label{eq:kyber_proof}
\end{equation}
Since $t = As + e$ (Eq.~\eqref{eq:mlwe}), the substitution yields $t^T r = s^T A^T r + e^T r$. The noise term $\nu = e^T r + e_2 - s^T e_1$ is a sum of products of small polynomials, bounded by $|\nu| < q/4$. The Compress rounding function (Eq.~\eqref{eq:kyber_compress})
\begin{equation}
m_{\text{final}} = \text{Compress}_q(m'_{poly}, 1) = \left\lfloor \frac{2}{q} m'_{poly} \right\rceil \!\!\bmod 2
\label{eq:kyber_compress}
\end{equation}
maps values near $0$ to bit $0$ and values near $q/2$ to bit $1$, correctly recovering $m$ with overwhelming probability. Verification confirms $ct' = ct$, yielding $SS_k = \text{KDF}(m, H(ct))$. $\hfill$

\textbf{Time Complexity.} The dominant operation is polynomial multiplication. With the Number-Theoretic Transform (NTT), each multiplication achieves $\mathcal{O}(n \log n)$ for $n = 256$. All three Kyber operations share the same asymptotic cost given by Eq.~\eqref{eq:kyber_complexity}:
\begin{equation}
\begin{split}
T_{\text{KEM.Gen}} &= T_{\text{KEM.Encap}} = T_{\text{KEM.Decap}} \\
                   &= \mathcal{O}(k^2 n \log n)
\end{split}
\label{eq:kyber_complexity}
\end{equation}
Since $k = 3$ and $n = 256$ are fixed, each algorithm effectively runs in constant time $\mathcal{O}(1)$.

\begin{figure}[!t]
\hrulefill
\vspace{-0.2cm}
\begin{center}
\textbf{Algorithm 4:} CRYSTALS-Kyber KEM (KeyGen, Encapsulation, Decapsulation)
\end{center}
\vspace{-0.2cm}
\hrulefill
\begin{enumerate}
    \item[] \textbf{--- KEM.KeyGen() ---}
    \item \textbf{Input:} Security parameter $k$, Ring dimension $n = 256$, Modulus $q = 3329$
    \item \textbf{Output:} Public key $pk = (A, t)$, Secret key $sk = s$
    \item $d \xleftarrow{\$} \{0,1\}^{256}$; \quad $(\rho, \sigma) \leftarrow \text{SHA3-512}(d)$
    \item \textbf{for} $i = 0$ \textbf{to} $k-1$, $j = 0$ \textbf{to} $k-1$ \textbf{do}
    \item \quad $A_{i,j} \leftarrow \text{Parse}(\text{SHAKE-128}(\rho \,\|\, i \,\|\, j))$ \quad \textit{// Rejection sampling in $R_q$}
    \item \textbf{end for}
    \item $(s, e) \leftarrow \text{CBD}_{\eta}(\sigma)$ \quad \textit{// Centered Binomial Distribution}
    \item $\hat{s} \leftarrow \text{NTT}(s)$; \quad $\hat{e} \leftarrow \text{NTT}(e)$
    \item $\hat{t} \leftarrow \hat{A} \circ \hat{s} + \hat{e}$ \quad \textit{// Pointwise multiply in NTT domain}
    \item \textbf{return} $pk \leftarrow (\hat{t}, \rho)$, $sk \leftarrow \hat{s}$
    \item[]
    \item[] \textbf{--- KEM.Encap($pk$) ---}
    \item \textbf{Input:} Public key $pk = (\hat{t}, \rho)$
    \item \textbf{Output:} Ciphertext $ct$, Shared secret $SS_k$
    \item $m \xleftarrow{\$} \{0,1\}^{256}$
    \item $(\bar{K}, r_{coins}) \leftarrow G(m \,\|\, H(pk))$
    \item Sample $(r, e_1, e_2)$ from $r_{coins}$ via CBD
    \item $u \leftarrow \text{NTT}^{-1}(\hat{A}^T \circ \text{NTT}(r)) + e_1$
    \item $v \leftarrow \text{NTT}^{-1}(\hat{t}^T \circ \text{NTT}(r)) + e_2 + \text{Encode}(m)$
    \item $ct \leftarrow (\text{Compress}_q(u, d_u), \text{Compress}_q(v, d_v))$
    \item $SS_k \leftarrow \text{KDF}(\bar{K} \,\|\, H(ct))$
    \item \textbf{return} $(ct, SS_k)$
    \item[]
    \item[] \textbf{--- KEM.Decap($sk, ct$) ---}
    \item \textbf{Input:} Secret key $sk$, Ciphertext $ct$
    \item \textbf{Output:} Shared secret $SS_k$
    \item $(u', v') \leftarrow \text{Decompress}(ct)$
    \item $m' \leftarrow \text{Decode}(v' - \text{NTT}^{-1}(\hat{s}^T \circ \text{NTT}(u')))$
    \item Re-encrypt: $ct' \leftarrow \text{KEM.Encap}_{det}(pk, m')$
    \item \textbf{if} $ct' = ct$ \textbf{then} $SS_k \leftarrow \text{KDF}(\bar{K} \,\|\, H(ct))$
    \item \textbf{else} $SS_k \leftarrow \text{KDF}(z \,\|\, H(ct))$ \quad \textit{// Implicit rejection}
    \item \textbf{return} $SS_k$
\end{enumerate}
\hrulefill
\end{figure}

\subsection{Protocol Lifecycle and Key Derivation Architecture}

The PQBFL protocol operates through four synchronized phases, each with formally defined cryptographic operations and blockchain interactions.

\subsubsection{Phase 1: On-Chain Identity Registration}



\textbf{Server Registration.} The aggregator server generates three independent key pairs:
\begin{equation}
\begin{split}
(kpk_b, ksk_b) &\leftarrow \text{Kyber.KeyGen()}, \\
(epk_b, esk_b) &\leftarrow \text{ECDH.Gen()}, \\
(pk_b, sk_b) &\leftarrow \text{ECDSA.Gen()}
\end{split}
\label{eq:server_keygen}
\end{equation}
A cryptographic commitment is constructed and published to the blockchain via the smart contract:
\begin{equation}    
\begin{split}
h_{keys} &\leftarrow \text{SHA3-256}(kpk_b \,\|\, epk_b), \\
h_{M0} &\leftarrow \text{SHA3-256}(M^0)
\end{split}
\label{eq:server_hash}
\end{equation}
\begin{equation}
\begin{split}
Tx_r &\leftarrow \{h_{keys}, n, h_{M0}, id_p\}, \\
BC &\leftarrow \text{RegisterProject}(Tx_r)
\end{split}
\label{eq:server_register}
\end{equation}

\textbf{Participant Registration.} Each healthcare participant generates ECDH and ECDSA pairs, computes $h_{epk_a} \leftarrow H(epk_a)$, and registers via:
\begin{equation}
\begin{split}
Tx_r' &\leftarrow \{h_{epk_a}, id_p\}, \\
BC &\leftarrow \text{RegisterClient}(Tx_r')
\end{split}
\label{eq:client_register}
\end{equation}

\begin{table}[H]
\centering
\caption{Phase 1 Computational Complexity per Entity}
\label{tab:phase1_complexity}
\begin{tabular}{l l l}
\toprule
\textbf{Operation} & \textbf{Complexity} & \textbf{Notes} \\
\midrule
Kyber.KeyGen() & $\mathcal{O}(k^2 n \log n)$ & NTT-accelerated \\
ECDH.Gen() & $\mathcal{O}(b^3)$ & One scalar mult. \\
ECDSA.Gen() & $\mathcal{O}(b^3)$ & One scalar mult. \\
SHA3-256 Hashing & $\mathcal{O}(L)$ & Linear in input \\
\bottomrule
\end{tabular}
\end{table}

\subsubsection{Phase 2: Hybrid Session Establishment}


The server transmits its public keys off-chain, signed with its ECDSA key:
\begin{equation}
\begin{split}
msg_{keys} &\leftarrow \text{Enc}(kpk_b \,\|\, epk_b), \\
\sigma_b &\leftarrow \text{ECDSA.Sign}(sk_b, msg_{keys})
\end{split}
\label{eq:session_send}
\end{equation}
Upon receipt, the participant verifies $\text{ECDSA.Verify}(pk_b, msg_{keys}, \sigma_b) \stackrel{?}{=} 1$ and performs the dual key agreement:
\begin{equation}
\begin{split}
S_{ecdh} &\leftarrow \text{ECDH.Derive}(esk_a, epk_b) \\
         &= esk_a \times epk_b
\end{split}
\label{eq:session_ecdh}
\end{equation}
\begin{equation}
(ct, S_{kem}) \leftarrow \text{Kyber.Encap}(kpk_b)
\label{eq:session_kem}
\end{equation}
The hybrid root key is derived through a two-stage HKDF-SHA256 extraction chain (Eq.~\eqref{eq:kdf_a_step1}--\eqref{eq:root_key}) to ensure crypto-agility:
\begin{equation}
prk_1 = \text{HMAC-SHA256}(\sigma,\ S_{kem})
\label{eq:kdf_a_step1}
\end{equation}
\begin{equation}
prk_2 = \text{HMAC-SHA256}(prk_1,\ S_{ecdh})
\label{eq:kdf_a_step2}
\end{equation}
\begin{equation}
RK_j = \text{HKDF-Expand}(prk_2,\ \texttt{"pqbfl:RK"},\ 32)
\label{eq:root_key}
\end{equation}
where $\sigma \in \{0,1\}^{256}$ is a session salt and the domain separation string prevents key confusion. The two-stage construction ensures that $RK_j$ (Eq.~\eqref{eq:root_key}) computationally commits to both primitives: if Kyber is broken by a quantum adversary, security holds under ECDH via Eq.~\eqref{eq:session_ecdh}, and vice versa.

\textbf{HKDF Mathematical Foundation.} HKDF is built on HMAC, defined as:
\begin{equation}
\begin{split}
\text{HMAC}(K, m) = H\big(&(K \oplus \text{opad}) \\
&\,\|\, H((K \oplus \text{ipad}) \,\|\, m)\big)
\end{split}
\label{eq:hmac_def}
\end{equation}
The Expand phase (Eq.~\eqref{eq:hkdf_expand}) iteratively generates output key material (OKM) using the HMAC construction of Eq.~\eqref{eq:hmac_def}:
\begin{equation}
\begin{split}
T(i) &= \text{HMAC}(\text{PRK},\ T(i\!-\!1) \,\|\, \text{label} \,\|\, 0x0i), \\
\text{OKM} &= T(1) \,\|\, T(2) \,\|\, \cdots
\end{split}
\label{eq:hkdf_expand}
\end{equation}
The time complexity is $\mathcal{O}(L_{\text{IKM}} + L)$, constant $\mathcal{O}(1)$ for fixed 32-byte inputs.

\begin{figure}[!t]
\hrulefill
\vspace{-0.2cm}
\begin{center}
\textbf{Algorithm 5:} Hybrid Session Establishment ($\text{KDF}_A$)
\end{center}
\vspace{-0.2cm}
\hrulefill
\begin{enumerate}
    \item \textbf{Input:} Server keys $(kpk_b, epk_b, sk_b)$, Participant keys $(esk_a, sk_a)$
    \item \textbf{Output:} Root key $RK_j$, Initial chain key $CK_{0,j}$
    \item[]
    \item[] \textbf{--- Server Side (Publish) ---}
    \item $msg_{keys} \leftarrow kpk_b \,\|\, epk_b$
    \item $\sigma_b \leftarrow \text{ECDSA.Sign}(sk_b, msg_{keys})$
    \item Transmit $\langle msg_{keys}, \sigma_b \rangle$ off-chain
    \item[]
    \item[] \textbf{--- Participant Side (Derive) ---}
    \item \textbf{Verify:} $\text{ECDSA.Verify}(pk_b, msg_{keys}, \sigma_b) \stackrel{?}{=} 1$
    \item \textbf{if} verification fails \textbf{then abort}
    \item $H_{recv} \leftarrow \text{SHA3-256}(kpk_b \,\|\, epk_b)$
    \item \textbf{Assert:} $H_{recv} = h_{keys}$ from on-chain $Tx_r$
    \item[]
    \item[] \textbf{--- Classical Agreement ---}
    \item $S_{ecdh} \leftarrow esk_a \times epk_b$ \quad \textit{// ECDH scalar mult.}
    \item[]
    \item[] \textbf{--- Post-Quantum Agreement ---}
    \item $(ct, S_{kem}) \leftarrow \text{Kyber.Encap}(kpk_b)$
    \item[]
    \item[] \textbf{--- Two-Stage Hybrid Extraction ---}
    \item $\sigma \xleftarrow{\$} \{0,1\}^{256}$ \quad \textit{// Random salt}
    \item $prk_1 \leftarrow \text{HMAC}(\sigma, S_{kem})$ \quad \textit{// Extract PQ secret}
    \item $prk_2 \leftarrow \text{HMAC}(prk_1, S_{ecdh})$ \quad \textit{// Extract classical secret}
    \item $RK_j \leftarrow \text{HKDF-Expand}(prk_2, \texttt{"pqbfl:RK"}, 32)$
    \item $CK_{0,j} \leftarrow \text{HMAC}(RK_j, \texttt{"pqbfl:CK0"})$
    \item \textbf{SecureErase}($S_{ecdh}, S_{kem}, prk_1, prk_2$)
    \item \textbf{return} $(RK_j, CK_{0,j}, ct, epk_a)$
\end{enumerate}
\hrulefill
\end{figure}

\subsubsection{Phase 3: Iterative Secure Model Training Loop}


During each FL training round $r$, the symmetric ratchet advances the key state. The server publishes a Task transaction containing the global model hash:
\begin{equation}
\begin{split}
Tx_p &\leftarrow \{r, H(Inf_b^r), id_p, id_t, D_t\}, \\
BC &\leftarrow \text{PublishTask}(Tx_p)
\end{split}
\label{eq:publish_task}
\end{equation}

Given the initial chain key $CK_{0,j} = \text{HMAC-SHA256}(RK_j, \texttt{"pqbfl:CK0"})$, each round $i$ advances via two domain-separated HMAC branches (Eq.~\eqref{eq:chain_key} and Eq.~\eqref{eq:model_key}):
\begin{equation}
CK_{i+1,j} = \text{HMAC}(CK_{i,j},\ \texttt{"pqbfl:CK"})
\label{eq:chain_key}
\end{equation}
\begin{equation}
K_{i,j} = \text{HMAC}(CK_{i,j},\ \texttt{"pqbfl:MK"})
\label{eq:model_key}
\end{equation}

\begin{figure*}
\centering
% Generate from Mermaid: mermaid_diagrams.md #6
\includegraphics[width=1.98\columnwidth]{fig3 (2).png}
\caption{Symmetric ratcheting key derivation tree ($\text{KDF}_S$). Each Chain Key $CK_i$ branches into the next chain key $CK_{i+1}$ and a disposable Model Key $K_i$ via domain-separated HMAC operations. The one-way structure ensures PRF-based forward secrecy: $K_{i-1}$ cannot be recovered from any $K_i$ or later state.}
\label{fig:kdf_tree}
\end{figure*}

After derivation per Eq.~\eqref{eq:chain_key}, $CK_{i,j}$ is immediately erased from memory. This structurally embeds Perfect Forward Secrecy (PFS): an adversary who obtains $K_{i,j}$ (Eq.~\eqref{eq:model_key}) or any future state cannot invert the HMAC to recover $CK_{i-1,j}$ or prior model keys. Domain labels \texttt{"pqbfl:CK"} and \texttt{"pqbfl:MK"} enforce PRF independence.

\textbf{Authenticated Encryption.} All federated model payloads are encrypted off-chain using AES-256-GCM:
\begin{equation}
\begin{split}
\eta &\xleftarrow{\$} \{0,1\}^{96}, \\
C_r &= \eta \;\|\; \text{AES-GCM.Enc}(K_{i,j}, \eta, \Delta_r, \mathcal{A})
\end{split}
\label{eq:aead_enc}
\end{equation}
where $\Delta_r$ is the local gradient update and $\mathcal{A} = \texttt{"pqbfl:}\langle\textit{dir}\rangle\texttt{:}r\texttt{"}$ is associated data (Eq.~\eqref{eq:aead_enc}). The GCM authentication uses polynomial hashing over Galois Field $GF(2^{128})$; the hash key and authentication tag are defined in Eq.~\eqref{eq:gcm_tag}:
\begin{equation}
\begin{split}
H_{auth} &= \text{AES}_K(0^{128}), \\
\text{Tag} &= \text{AES}_K(\text{Nonce}) \oplus \text{GHASH}(H_{auth}, \mathcal{A}, C)
\end{split}
\label{eq:gcm_tag}
\end{equation}
where the GHASH function (Eq.~\eqref{eq:ghash}) computes a polynomial hash:
\begin{equation}
\begin{split}
\text{GHASH} = &\sum_{i} A_i \cdot H_{auth}^{n-i} \\
              +&\sum_{j} C_j \cdot H_{auth}^{m-j} \pmod{P}
\end{split}
\label{eq:ghash}
\end{equation}
with $P$ being a fixed irreducible polynomial in $GF(2^{128})$. The time complexity is $\mathcal{O}(N_A + N_M) = \mathcal{O}(L_A + L_M)$, linear in data length.

\textbf{Training, Upload, and Aggregation.} The participant trains locally producing $m^r \leftarrow \text{Train}(M^{r-1}, \text{Data}_{local})$, encrypts and signs:
\begin{equation}
\begin{split}
msg_a &\leftarrow \text{Enc}((Inf_a^r \,\|\, Tx_u),\ K_{i,j}), \\
\sigma_a &\leftarrow \text{ECDSA.Sign}(sk_a, msg_a)
\end{split}
\label{eq:participant_send}
\end{equation}
The server verifies, decrypts, checks integrity against the on-chain hash, and aggregates the verified local models via Federated Averaging (Eq.~\eqref{eq:fedavg}):
\begin{equation}
M^r = \sum_{i=1}^{N} \frac{w_i}{N} m_i^r
\label{eq:fedavg}
\end{equation}

\begin{figure}[!t]
\hrulefill
\vspace{-0.2cm}
\begin{center}
\textbf{Algorithm 6:} End-to-End Secure Federated Learning Round
\end{center}
\vspace{-0.2cm}
\hrulefill
\begin{enumerate}
    \item \textbf{Input:} Global model $M^{r-1}$, Round $r$, Chain key $CK_{i-1,j}$, Participant set $\mathcal{P}$
    \item \textbf{Output:} Updated global model $M^r$, Updated chain key $CK_{i,j}$
    \item[]
    \item[] \textbf{--- Server: Publish Task ---}
    \item $(CK_{i,j}, K_{i,j}) \leftarrow \text{KDF}_S(CK_{i-1,j})$; \quad \textbf{SecureErase}($CK_{i-1,j}$)
    \item $Inf_b^r \leftarrow \langle r, M^{r-1}, id_p, id_t, D_t \rangle$
    \item $Tx_p \leftarrow \{r, H(Inf_b^r), id_p, id_t, D_t\}$
    \item $BC \leftarrow \text{PublishTask}(Tx_p)$ \quad \textit{// On-chain commitment}
    \item $\eta \xleftarrow{\$} \{0,1\}^{96}$
    \item $msg_b \leftarrow \eta \,\|\, \text{AES-GCM.Enc}(K_{i,j}, \eta, Inf_b^r \,\|\, Tx_p, \mathcal{A}_s)$
    \item $\sigma_b \leftarrow \text{ECDSA.Sign}(sk_b, msg_b)$
    \item Broadcast $\langle msg_b, \sigma_b \rangle$ to all $P_k \in \mathcal{P}$
    \item[]
    \item[] \textbf{--- Participant $P_k$: Train and Upload ---}
    \item \textbf{Verify:} $\text{ECDSA.Verify}(pk_b, msg_b, \sigma_b) \stackrel{?}{=} 1$
    \item $(CK_{i,j}, K_{i,j}) \leftarrow \text{KDF}_S(CK_{i-1,j})$
    \item $(Inf_b^r \,\|\, Tx_p) \leftarrow \text{AES-GCM.Dec}(K_{i,j}, msg_b)$
    \item \textbf{Assert:} $H(Inf_b^r) = h$ from on-chain $Tx_p$
    \item $m_k^r \leftarrow \text{Train}(M^{r-1}, \text{Data}_{local}^k)$
    \item $Tx_u \leftarrow \{r, H(m_k^r), id_p, id_t\}$
    \item $BC \leftarrow \text{UpdateModel}(Tx_u)$ \quad \textit{// On-chain proof of contribution}
    \item $msg_a \leftarrow \text{AES-GCM.Enc}(K_{i,j}, \eta', m_k^r \,\|\, Tx_u, \mathcal{A}_p)$
    \item $\sigma_a \leftarrow \text{ECDSA.Sign}(sk_a, msg_a)$
    \item Transmit $\langle msg_a, \sigma_a \rangle$ to server
    \item[]
    \item[] \textbf{--- Server: Verify and Aggregate ---}
    \item \textbf{for each} $P_k \in \mathcal{P}$ \textbf{do}
    \item \quad $\text{ECDSA.Verify}(pk_{a_k}, msg_{a_k}, \sigma_{a_k})$
    \item \quad $m_k^r \leftarrow \text{AES-GCM.Dec}(K_{i,j}, msg_{a_k})$
    \item \quad \textbf{Assert:} $\text{hmac.compare\_digest}(H(m_k^r), h_{\text{on-chain}})$
    \item \textbf{end for}
    \item $M^r \leftarrow \sum_{k=1}^{N} \frac{w_k}{N} m_k^r$ \quad \textit{// FedAvg aggregation}
    \item $Tx_f \leftarrow \{r, sc, H(M^r), id_p, id_t\}$
    \item $BC \leftarrow \text{FeedbackModel}(Tx_f)$
    \item \textbf{return} $(M^r, CK_{i,j})$
\end{enumerate}
\hrulefill
\end{figure}

\begin{figure}[!t]
\hrulefill
\vspace{-0.2cm}
\begin{center}
\textbf{Algorithm 7:} Symmetric Ratchet Key Derivation ($KDF_S$)
\end{center}
\vspace{-0.2cm}
\hrulefill
\begin{enumerate}
    \item \textbf{Input:} Root Key $RK_j$, Round Index $i$, Model Payload $\Delta_r$
    \item \textbf{Output:} Ciphertext $C_r$, Authentication Tag $\tau$
    \item Initialize chain: $CK_{0,j} \leftarrow \text{HMAC}(RK_j, \texttt{"pqbfl:CK0"})$
    \item \textbf{for} $i = 1$ \textbf{to} $L_j$ \textbf{do}
    \item \quad $CK_{i,j} \leftarrow \text{HMAC}(CK_{i-1,j}, \texttt{"pqbfl:CK"})$
    \item \quad $K_{i,j} \leftarrow \text{HMAC}(CK_{i-1,j}, \texttt{"pqbfl:MK"})$
    \item \quad \textbf{SecureErase}($CK_{i-1,j}$)
    \item \quad $\eta \xleftarrow{\$} \{0,1\}^{96}$
    \item \quad $C_r \leftarrow \eta \,\|\, \text{AES-GCM.Enc}(K_{i,j}, \eta, \Delta_r, \mathcal{A})$
    \item \quad $\sigma \leftarrow \text{ECDSA.Sign}(sk, C_r)$
    \item \quad \textbf{Transmit} $\langle C_r, \sigma \rangle$ off-chain
    \item \textbf{end for}
\end{enumerate}
\hrulefill
\end{figure}
\newpage

\subsection{Threat-Adaptive Ratcheting Framework}
\label{sec:contribution_adaptive}

\textbf{Our core contribution} eliminates the static ratcheting bottleneck by introducing a dynamic threat evaluation loop over live blockchain telemetry. The rigid per-session threshold $L_j$ of the base protocol is overridden at runtime by a continuously updated adaptive policy.

\subsubsection{ThreatMonitor: Composite Threat Level Computation}

\begin{figure*}[!t]
\centering
% Generate from Mermaid: mermaid_diagrams.md #7
\includegraphics[width=\textwidth]{fig4 (1).png}
\caption{\texttt{ThreatMonitor} event aggregation pipeline. Five cryptographically-derived signal categories (weighted by severity and recency via exponential decay $d_k = 2^{-\Delta\tau/\lambda}$) are combined into a normalized composite threat score $t \in [0,1]$, which drives the \texttt{AdaptiveRatchetPolicy}.}
\label{fig:threatmonitor}
\end{figure*}

The \texttt{ThreatMonitor} module continuously aggregates security signals from blockchain transaction traces. Let $\mathcal{E} = \{(e_i, s_i, \tau_i)\}_{i=1}^{N}$ denote recorded security events, where $e_i$ is the event type, $s_i \in [0, 1]$ is caller-supplied severity, and $\tau_i$ is the UNIX epoch timestamp. The monitor considers events within a sliding window of width $W$ seconds (default $W = 300$\,s). For each active event, an exponential decay factor is computed:
\begin{equation}
d_i = 2^{-A_i / \lambda}, \quad A_i = \tau_{now} - \tau_i
\label{eq:decay_factor}
\end{equation}
where $\lambda > 0$ is the half-life constant (default $\lambda = 120$\,s) and the decay factor is given by Eq.~\eqref{eq:decay_factor}. The composite threat level is the exponentially decayed, type-weighted average (Eq.~\eqref{eq:threat_level}):
\begin{equation}
t = \frac{\displaystyle\sum_{i:\, A_i \leq W} d_i \cdot w_{e_i} \cdot s_i}{\displaystyle\sum_{i:\, A_i \leq W} d_i \cdot w_{e_i}}
\label{eq:threat_level}
\end{equation}
where $w_{e_i}$ is the category weight and $t \in [0, 1]$. The five event categories and weights are: \texttt{sig\_verification\_failed} ($w = 1.0$), \texttt{hash\_mismatch} ($w = 0.9$), \texttt{reputation\_drop} ($w = 0.6$), \texttt{timing\_anomaly} ($w = 0.4$), and \texttt{stale\_ratchet} ($w = 0.3$).

\textbf{Example Calculation.} Assume $\lambda = 10$\,s. Within a 5-second interval, 3 signature failures occur ($\Delta\tau = 5, 2, 0$\,s, $sev = 1.0$) plus a timing anomaly ($\Delta\tau = 1$\,s, $sev = 0.5$):
\begin{equation}
\begin{split}
\text{decay}_{e_1} &= 2^{-5/10} + 2^{-2/10} + 2^{0} \\
                   &\approx 0.707 + 0.871 + 1.0 = 2.578
\end{split}
\label{eq:decay_example}
\end{equation}
Raw score: $(2.578 \times 1.0 \times 1.0) + (0.933 \times 0.4 \times 0.5) = 2.765$. If max aggregate is 5.0, then $t = 0.553$.

\subsubsection{AdaptiveRatchetPolicy: Threat-to-$L_j$ Mapping}
\begin{figure}[!t]
\centering
\includegraphics[width=0.5\textwidth]{fig5 (2).png}
\caption{\texttt{AdaptiveRatchetPolicy} decision graph. The raw threat score $t$ is mapped through the power-curve function $L_j(t)=\lfloor L_{max}-(L_{max}-L_{min})\cdot t^\gamma \rceil$, bounded to $[L_{min},L_{max}]$, and gated by a cooldown check before being published as the active symmetric ratcheting window $L_j^{adaptive}$.}
\label{fig:adaptive_policy}
\end{figure}

The policy maps $t$ to a concrete symmetric ratcheting window $L_j$ via the power-curve function in Eq.~\eqref{eq:adaptive_policy}:
\begin{equation}
L_j(t) = \left\lfloor L_{max} - (L_{max} - L_{min}) \cdot t^{\gamma} \right\rceil
\label{eq:adaptive_policy}
\end{equation}
with $L_{min} = 2$, $L_{max} = 20$, and $\gamma = 2.0$ (quadratic sensitivity). This mapping provides three formally verifiable properties:

\textbf{Property 1 (Monotone Decrease):} $\forall\, t_1 < t_2 \in [0,1]:\; L_j(t_1) \geq L_j(t_2)$. Higher threat always produces tighter rotation.

\textbf{Property 2 (Non-Linear Noise Suppression):} For $\gamma > 1$, the marginal reduction in $L_j$ per unit increase in $t$ is sublinear near $t = 0$, meaning transient events do not trigger expensive KEM re-keying. At high threat ($t \to 1$), $L_j \to L_{min}$ rapidly.

\textbf{Property 3 (Boundary Safety):} $L_j \in [L_{min}, L_{max}]$ ensures the system never degenerates into per-round asymmetric re-keying (prohibitively expensive) nor extends indefinitely without re-keying (collapsing PCS).

\textbf{Example.} Using $t = 0.553$: $L_j(0.553) = \lfloor 20 - 18 \times (0.553)^2 \rfloor = \lfloor 20 - 5.508 \rfloor = 14$.

\begin{figure}[!t]
\centering
\includegraphics[width=\columnwidth]{graph2_lj_vs_threat.png}
\caption{Adaptive ratcheting window $L_j$ as a function of composite threat level $t \in [0,1]$ (Eq.~\eqref{eq:adaptive_policy}), with $\gamma = 2.0$, $L_{min} = 2$, $L_{max} = 20$. The green zone indicates where the adaptive protocol outperforms the fixed baseline ($L_j > 10$); the blue zone shows where the adaptive policy is more conservative than the baseline under high threat. The current zero-threat operating point ($L_j = 20$) is marked by the green dot.}
\label{fig:lj_vs_threat_curve}
\end{figure}

A cooldown gate (Eq.~\eqref{eq:cooldown_gate}) prevents oscillation:
\begin{equation}
L_j^{(r)} \!=\! \begin{cases} L_j(t) & \!\text{if } |r\!-\!r_{last}| \!\geq\! C_{cd} \\ & \text{and } L_j(t) \!\neq\! L_j^{(r-1)} \\ L_j^{(r-1)} & \!\text{otherwise} \end{cases}
\label{eq:cooldown_gate}
\end{equation}

\subsubsection{Conservative Dual-Threshold Ratchet Trigger}


The protocol triggers a full asymmetric ratchet when:
\begin{equation}
\begin{split}
L_j^{eff} &= \min\!\left(L_j^{adaptive},\ L_j(t)\right), \\
\text{RatchetNow} &= \mathbb{1}\!\left[i \geq L_j^{eff}\right]
\end{split}
\label{eq:dual_threshold}
\end{equation}
This conservative minimum (Eq.~\eqref{eq:dual_threshold}) ensures that a sudden threat spike immediately tightens the remaining window, even if the cooldown gate of Eq.~\eqref{eq:cooldown_gate} has not yet updated.

\begin{figure}[!t]
\hrulefill
\vspace{-0.2cm}
\begin{center}
\textbf{Algorithm 8:} Threat-Adaptive Ratchet Policy Evaluation
\end{center}
\vspace{-0.2cm}
\hrulefill
\begin{enumerate}
    \item \textbf{Input:} Event Stream $\mathcal{E}$, Current round $i$, Session index $j$, Parameters $(\lambda, W, \gamma, L_{min}, L_{max}, C_{cd})$
    \item \textbf{Output:} Updated $L_j$, Ratchet trigger decision
    \item $\mathcal{E}_{active} \leftarrow \{(e_k, s_k, \tau_k) \in \mathcal{E} \mid \tau_{now} - \tau_k \leq W\}$
    \item \textbf{for each} $(e_k, s_k, \tau_k) \in \mathcal{E}_{active}$ \textbf{do}
    \item \quad $d_k \leftarrow 2^{-(\tau_{now} - \tau_k)/\lambda}$
    \item \quad $\text{num} \leftarrow \text{num} + d_k \cdot w_{e_k} \cdot s_k$; \quad $\text{den} \leftarrow \text{den} + d_k \cdot w_{e_k}$
    \item \textbf{end for}
    \item $t \leftarrow \text{num} / \text{den}$ \quad (clamp to $[0, 1]$)
    \item $L_j^{new} \leftarrow \lfloor L_{max} - (L_{max} - L_{min}) \cdot t^{\gamma} \rceil$
    \item \textbf{if} $|i - r_{last}| \geq C_{cd}$ \textbf{and} $L_j^{new} \neq L_j^{current}$ \textbf{then}
    \item \quad $L_j^{current} \leftarrow L_j^{new}$; \quad $r_{last} \leftarrow i$
    \item \quad Log $(i, L_j^{old} \to L_j^{new}, t)$ to on-chain audit
    \item \textbf{end if}
    \item $L_j^{eff} \leftarrow \min(L_j^{current}, L_j^{new})$
    \item \textbf{if} $i \geq L_j^{eff}$ \textbf{then}
    \item \quad \textbf{Trigger Phase 4: Asymmetric Ratchet}
    \item \quad Generate fresh $(kpk'_b, ksk'_b) \leftarrow \text{Kyber.Gen()}$, $(epk'_b, esk'_b) \leftarrow \text{ECDH.Gen()}$
    \item \quad Derive $RK_{j+1}$ via Eq.~\eqref{eq:root_key} with new salt $\sigma' \xleftarrow{\$} \{0,1\}^{256}$
    \item \quad Reset $i \leftarrow 0$, $j \leftarrow j+1$
    \item \textbf{end if}
\end{enumerate}
\hrulefill
\end{figure}

\subsubsection{Phase 4: Asymmetric Ratchet Execution}
When $\text{RatchetNow}=1$, the server generates fresh Kyber and ECDH key pairs, piggybacks them onto the PublishTask message encrypted with the last key $K_{i,j}$:
\begin{equation}
\begin{split}
Inf_b^r \leftarrow \langle r, kpk'_b,& epk'_b, \\
                  & M^{r-1}, id_p, id_t \rangle
\end{split}
\label{eq:ratchet_info}
\end{equation}
\begin{equation}
\begin{split}
msg_b &\leftarrow \text{Enc}((Inf_b^r \,\|\, Tx_p),\ K_{i,j}), \\
\sigma_b &\leftarrow \text{ECDSA.Sign}(sk_b, msg_b)
\end{split}
\label{eq:ratchet_send}
\end{equation}
The participant decrypts, extracts new keys, generates new ephemeral ECDH pair $(epk'_a, esk'_a)$, computes new dual secrets (Eq.~\eqref{eq:ratchet_secrets}), and derives the new session root (Eq.~\eqref{eq:ratchet_derive}):
\begin{equation}
\begin{split}
SS_e &\leftarrow \text{ECDH.Derive}(esk'_a, epk'_b), \\
(ct', SS_k) &\leftarrow \text{Kyber.Encap}(kpk'_b)
\end{split}
\label{eq:ratchet_secrets}
\end{equation}
\begin{equation}
\begin{split}
RK_{j+1} &\leftarrow \text{KDF}_A(SS_e \,\|\, SS_k,\ \sigma'), \\
(CK_{1,j+1}, K_{1,j+1}) &\leftarrow \text{KDF}_S(RK_{j+1})
\end{split}
\label{eq:ratchet_derive}
\end{equation}
The participant encrypts her model response with the \textit{first key of the new chain} $K_{1,j+1}$, completing the ratchet. The server decapsulates $SS_k \leftarrow \text{Kyber.Decap}(ksk'_b, ct')$, independently derives the identical $RK_{j+1}$ and $K_{1,j+1}$, and decrypts.

\begin{figure}[!t]
\hrulefill
\vspace{-0.2cm}
\begin{center}
\textbf{Algorithm 9:} Asymmetric Ratchet Execution (Phase 4 --- Server and Participant)
\end{center}
\vspace{-0.2cm}
\hrulefill
\begin{enumerate}
    \item \textbf{Input:} Current keys $(ksk_b, esk_b, sk_b)$, Last model key $K_{i,j}$, Round $r$
    \item \textbf{Output:} New root key $RK_{j+1}$, New chain $CK_{1,j+1}$
    \item[]
    \item[] \textbf{--- Server: Initiate Ratchet ---}
    \item $(kpk'_b, ksk'_b) \leftarrow \text{Kyber.KeyGen()}$
    \item $(epk'_b, esk'_b) \leftarrow \text{ECDH.Gen()}$
    \item $Inf_b^r \leftarrow \langle r, kpk'_b, epk'_b, M^{r-1}, id_p, id_t \rangle$
    \item $Tx_p \leftarrow \{r, H(Inf_b^r), H(kpk'_b \,\|\, epk'_b), id_p, id_t\}$
    \item $BC \leftarrow \text{PublishTask}(Tx_p)$ \quad \textit{// Commit new key hashes on-chain}
    \item $msg_b \leftarrow \text{AES-GCM.Enc}(K_{i,j}, \eta, Inf_b^r \,\|\, Tx_p, \mathcal{A})$ \quad \textit{// Encrypt with OLD key}
    \item $\sigma_b \leftarrow \text{ECDSA.Sign}(sk_b, msg_b)$
    \item Transmit $\langle msg_b, \sigma_b \rangle$
    \item[]
    \item[] \textbf{--- Participant: Respond to Ratchet ---}
    \item Verify $\sigma_b$; Decrypt $Inf_b^r$ using $K_{i,j}$
    \item \textbf{Assert:} $\text{hmac.compare\_digest}(H(kpk'_b \,\|\, epk'_b), h_{\text{on-chain}})$
    \item Extract $(kpk'_b, epk'_b)$ from $Inf_b^r$
    \item $(epk'_a, esk'_a) \leftarrow \text{ECDH.Gen()}$ \quad \textit{// Fresh ephemeral}
    \item $SS_e \leftarrow esk'_a \times epk'_b$ \quad  \textit{// New ECDH secret}
    \item $(ct', SS_k) \leftarrow \text{Kyber.Encap}(kpk'_b)$ \quad \textit{// New KEM secret}
    \item $\sigma' \xleftarrow{\$} \{0,1\}^{256}$ \quad \textit{// Fresh random salt}
    \item $RK_{j+1} \leftarrow \text{KDF}_A(SS_e \,\|\, SS_k, \sigma')$
    \item $(CK_{1,j+1}, K_{1,j+1}) \leftarrow \text{KDF}_S(RK_{j+1})$
    \item \textbf{SecureErase}($SS_e, SS_k, K_{i,j}, CK_{i,j}$)
    \item $m^r \leftarrow \text{Train}(M^{r-1}, \text{Data}_{local})$
    \item $msg_a \leftarrow \text{AES-GCM.Enc}(K_{1,j+1}, \eta', m^r \,\|\, ct' \,\|\, epk'_a, \mathcal{A}')$ \quad \textit{// Encrypt with NEW key}
    \item $\sigma_a \leftarrow \text{ECDSA.Sign}(sk_a, msg_a)$
    \item Transmit $\langle msg_a, \sigma_a \rangle$
    \item[]
    \item[] \textbf{--- Server: Complete Ratchet ---}
    \item Verify $\sigma_a$; Extract $(ct', epk'_a)$ from $Tx_u$
    \item $SS_e \leftarrow esk'_b \times epk'_a$; \quad $SS_k \leftarrow \text{Kyber.Decap}(ksk'_b, ct')$
    \item $RK_{j+1} \leftarrow \text{KDF}_A(SS_e \,\|\, SS_k, \sigma')$
    \item $(CK_{1,j+1}, K_{1,j+1}) \leftarrow \text{KDF}_S(RK_{j+1})$
    \item Decrypt $msg_a$ using $K_{1,j+1}$; Aggregate; Post $Tx_f$
    \item \textbf{SecureErase}($ksk_b, esk_b, SS_e, SS_k$) \quad \textit{// Old keys destroyed}
    \item Reset $i \leftarrow 0$, $j \leftarrow j+1$
    \item \textbf{return} $(RK_{j+1}, CK_{1,j+1})$
\end{enumerate}
\hrulefill
\end{figure}

\subsection{Security Proofs for the Adaptive Framework}

\textbf{Theorem 4 (Forward Secrecy Preservation).} The adaptive ratcheting mechanism preserves identical forward secrecy guarantees as the base protocol.

\textit{Proof.} Forward secrecy rests on the security of $\text{KDF}_S$. Let $\text{KDF}_S: \{0,1\}^{\lambda} \to \{0,1\}^{\lambda} \times \{0,1\}^{\lambda}$ be a one-way function mapping $CK_{i-1} \mapsto (CK_i, K_i)$. Because $\text{KDF}_S$ is modeled as a Pseudo-Random Function (PRF), reversing the chain to find $CK_{i-1}$ given $CK_i$ requires breaking the PRF assumption. The dynamic variation of $L_j$ strictly controls the quantity of iterations $i$, not the mathematical structure of $\text{KDF}_S$. Since $\forall\, i$, $\Pr[\text{Invert KDF}_S] \leq \epsilon_{\text{PRF}}$, altering $L_j$ dynamically preserves identical forward secrecy bounds. $\hfill$

\textbf{Theorem 5 (Enhanced Post-Compromise Security Window).} Under the adaptive scheme, the post-compromise vulnerability window is strictly reduced.

\textit{Proof.} Post-compromise security recovery occurs at each asymmetric ratchet deriving $RK_{j+1}$. The attacker's window of compromise is bounded by:
\begin{equation}
|K_{\text{compromise}}| = L_j - c, \quad \text{where } c \text{ is the round of compromise}
\label{eq:pcs_window}
\end{equation}
Under fixed ratcheting: $|K_{\text{compromise}}^{\text{fixed}}| = L_j^{\text{fixed}} - c$. Under adaptive ratcheting, active attack drives $t \to 1.0$, forcing $L_j^{\text{adaptive}} \to L_{min}$. Since $L_{min} \ll L_j^{\text{fixed}}$:
\begin{equation}
|K_{\text{compromise}}^{\text{adaptive}}| < |K_{\text{compromise}}^{\text{fixed}}|
\label{eq:pcs_comparison}
\end{equation}
The adversary advantage remains bounded by:
\begin{equation}
Adv_{\mathcal{A}}^{\text{PCS}} \leq \min\!\left(Adv_{\mathcal{A},\text{Kyber}}^{\text{IND-CCA}},\ Adv_{\mathcal{A},\text{ECDH}}^{\text{PRF-ODH}}\right)
\label{eq:adversary_advantage}
\end{equation}
but the frequency of PCS recovery truncates the utility lifespan of any compromised symmetric subset. $\hfill$

\subsection{Implementation Hardening Against Side-Channel Attacks}



\label{sec:contribution_hardening}

The base PQBFL codebase contained three implementation-level vulnerabilities exploitable independently of the algorithmic security proofs. We apply focused, overhead-minimal mitigations. (i)~All cryptographic comparisons are replaced with \texttt{hmac.compare\_digest()}, enforcing $T_{\text{compare}}(x,y) = \Theta(n)$ and eliminating byte-by-byte timing oracles. (ii)~AES-GCM nonces change from deterministic $\eta = H(\texttt{label}\|r)$ to $\eta \xleftarrow{\$} \{0,1\}^{96}$. (iii)~The fixed HKDF-Extract salt $\sigma = 0^{32}$ is replaced by $\sigma \xleftarrow{\$} \{0,1\}^{256}$ at each asymmetric ratchet boundary.

\textbf{Nonce Collision Probability.} Under a 96-bit true random nonce, the collision probability before key rotation is bounded by the Birthday Paradox:
\begin{equation}
P_c(N) = 1 - e^{-N^2 / 2 \cdot 2^{96}} \approx \frac{N^2}{2^{97}}
\label{eq:nonce_collision}
\end{equation}
Since the symmetric ratchet rotates $K_{i,j}$ per round, $N$ remains essentially 1, giving $P_c(1) \approx 2^{-97}$, rendering collision operationally impossible.

\textbf{Random Salt Search Space Inflation.} Randomizing the KDF salt inflates the precomputation search space:
\begin{equation}
S_{\text{dynamic}} = S_{\text{static}} \times 2^{256} = \mathcal{O}(2^{256})
\label{eq:salt_inflation}
\end{equation}
This guarantees computational infeasibility against rainbow table attacks while inducing zero execution latency penalty ($\mathcal{O}(1)$ derivation).

\begin{figure}[!t]
\hrulefill
\vspace{-0.2cm}
\begin{center}
\textbf{Algorithm 10:} Side-Channel Hardened Cryptographic Operations
\end{center}
\vspace{-0.2cm}
\hrulefill
\begin{enumerate}
    \item \textbf{Input:} Values $A, B$ to compare, Nonce mode, Salt mode
    \item \textbf{Output:} Hardened comparison result, Hardened nonce, Hardened salt
    \item[]
    \item[] \textbf{--- Constant-Time Comparison ---}
    \item \textbf{function} \textsc{SecureCompare}($A, B$)
    \item \quad $\text{result} \leftarrow 0$
    \item \quad \textbf{for} $i = 0$ \textbf{to} $|A| - 1$ \textbf{do}
    \item \quad \quad $\text{result} \leftarrow \text{result} \;|\; (A[i] \oplus B[i])$ \quad \textit{// Bitwise OR accumulation}
    \item \quad \textbf{end for} \quad \textit{// Always traverses full length $\Theta(n)$}
    \item \quad \textbf{return} $(\text{result} = 0)$
    \item[]
    \item[] \textbf{--- Randomized Nonce Generation ---}
    \item \textbf{function} \textsc{SecureNonce}()
    \item \quad $\eta \xleftarrow{\$} \{0,1\}^{96}$ \quad \textit{// CSPRNG, not H(label$\|$r)}
    \item \quad \textbf{Assert:} $\eta \notin \mathcal{N}_{\text{used}}$ \quad \textit{// Nonce uniqueness check}
    \item \quad $\mathcal{N}_{\text{used}} \leftarrow \mathcal{N}_{\text{used}} \cup \{\eta\}$
    \item \quad \textbf{return} $\eta$
    \item[]
    \item[] \textbf{--- Random KDF Salt ---}
    \item \textbf{function} \textsc{SecureSalt}()
    \item \quad $\sigma \xleftarrow{\$} \{0,1\}^{256}$ \quad \textit{// Not $0^{32}$ static}
    \item \quad \textbf{return} $\sigma$
    \item[]
    \item[] \textbf{--- Hardened Hash Verification ---}
    \item \textbf{function} \textsc{VerifyCommitment}($\text{data}, h_{\text{chain}}$)
    \item \quad $h_{\text{local}} \leftarrow \text{SHA3-256}(\text{data})$
    \item \quad \textbf{return} \textsc{SecureCompare}($h_{\text{local}}, h_{\text{chain}}$) \quad \textit{// No timing leak}
\end{enumerate}
\hrulefill
\end{figure}

\begin{table}[H]
\centering
\caption{Side-Channel Hardening Summary}
\label{tab:sc_hardening}
\begin{tabular}{l l l}
\toprule
\textbf{Vulnerability} & \textbf{Base Behavior} & \textbf{Fix} \\
\midrule
Timing oracle & \texttt{==} early-exit & \texttt{hmac.compare\_digest} \\
Nonce reuse & $\eta = H(\texttt{label}\|r)$ & $\eta \xleftarrow{\$}\{0,1\}^{96}$ \\
KDF precomputation & $\sigma = 0^{32}$ fixed & $\sigma \xleftarrow{\$}\{0,1\}^{256}$ \\
\bottomrule
\end{tabular}
\end{table}

\subsection{Comprehensive Complexity Summary}

\begin{table}[H]
\centering
\caption{End-to-End Per-Round Computational Complexity}
\label{tab:complexity_summary}
\begin{tabular}{l l l}
\toprule
\textbf{Operation} & \textbf{Complexity} & \textbf{Practical Time} \\
\midrule
ECDSA Sign/Verify & $\mathcal{O}(b^3)$ & $\sim 0.5$\,ms \\
ECDH Agreement & $\mathcal{O}(b^3)$ & $\sim 0.3$\,ms \\
Kyber KEM (all ops) & $\mathcal{O}(k^2 n \log n)$ & $\sim 3$--$5$\,ms \\
HMAC-SHA256 (KDF\_S) & $\mathcal{O}(1)$ & $\sim 1$\,$\mu$s \\
AES-256-GCM & $\mathcal{O}(L_M)$ & $\sim 10$\,$\mu$s/MB \\
ThreatMonitor Eval. & $\mathcal{O}(1)$ & $\sim 10$\,$\mu$s \\
AdaptivePolicy & $\mathcal{O}(1)$ & $\sim 1$\,$\mu$s \\
\bottomrule
\end{tabular}
\end{table}


\section{Experimental Results and Performance Evaluation}
To validate the Threat-Adaptive PQBFL framework, a full end-to-end prototype was deployed comprising 50 federated learning rounds across 10 participant clients connected to an Ethereum-compatible blockchain (Ganache, \texttt{http://127.0.0.1:8545}). The CRYSTALS-Kyber backend was configured as \texttt{ml\_kem\_512} (NIST FIPS~203, 128-bit post-quantum security). A total of 1,062 on-chain transactions were generated and logged across the experiment, spanning contract deployment, client registration, task publication, model updates, and feedback commitments. Table~\ref{tab:sim_params} summarizes all configuration parameters used throughout the evaluation.

\begin{table}[H]
\centering
\caption{Simulation and Protocol Configuration Parameters}
\label{tab:sim_params}
\begin{tabular}{l l l}
\toprule
\textbf{Parameter} & \textbf{Value} & \textbf{Notes} \\
\midrule
Federated Rounds ($R$)             & 50              & Full protocol lifecycle \\
Clients ($N$)                      & 10              & Participants $P_1$--$P_{10}$ \\
KEM Backend                        & \texttt{ml\_kem\_512}   & NIST FIPS 203, Level 1 \\
Blockchain                         & Ganache         & Local Ethereum testnet \\
$L_{min}$ (max security)           & 2               & Minimum ratchet window \\
$L_{max}$ (max efficiency)         & 20              & Maximum ratchet window \\
$L_{default}$ (initial)            & 10              & Starting value \\
Sensitivity exponent $\gamma$      & 0.50            & Power-curve shaping \\
Total on-chain transactions        & 1,062           & All phases combined \\
Initial model accuracy             & 60.50\%         & Round 0 \\
Final model accuracy               & 81.75\%         & Round 50 \\
Accuracy improvement               & $+21.25\%$      & Absolute gain \\
Avg.\ blockchain tx latency        & 8.00\,ms        & Per-transaction average \\
Min tx latency                     & 6.73\,ms        & Best-case per tx \\
Max tx latency                     & 36.48\,ms       & Cold-start contract deploy \\
Avg.\ crypto op latency            & $\sim$3--5\,ms  & AES-GCM, HMAC, KDF$_S$ \\
Security checks passed             & 10/10           & All hardening verified \\
\bottomrule
\end{tabular}
\end{table}

\subsection{Experimental Setup}
The participant nodes simulate a distributed hospital data consortium where patient records are processed locally without leaving the device. The global aggregation server runs the FedAvg aggregator and blockchain oracle on an Ethereum-compatible Ganache node. Smart contracts govern the immutable on-chain logging of \texttt{ThreatMonitor} composite scores, ECDSA-signed session key hashes, and gradient update commitments $H(m_k^r)$.

The adaptive ratcheting parameters were configured with $L_{min} = 2$, $L_{max} = 20$, $L_{default} = 10$, and sensitivity $\gamma = 0.50$. The \texttt{ThreatMonitor} applies exponential decay $d_k = 2^{-\Delta\tau/\lambda}$ to all logged security events, ensuring older events contribute negligibly as the system recovers. All cryptographic operations execute via constant-time C bindings through the \texttt{liboqs} Open Quantum Safe library.

\subsection{Testbed Setup and Threat Scenario}
To validate the adaptive ratcheting mechanism, ten participant nodes were configured to emulate a realistic adversarial scenario. All 10 sessions simultaneously accumulated \texttt{stale\_ratchet} events (severity weight $w=0.30$) as the symmetric ratchet window grew without a Phase~4 reset, triggering the \texttt{ThreatMonitor} to compute a peak composite threat score of $t = 0.879$ at round~3.

Fig.~\ref{fig:L_j_threat_combined} presents the dual-axis visualization of the adaptive ratcheting window $L_j$ and composite threat score $t$ over all 50 rounds, which is the primary empirical result of this work.

\begin{figure*}[!t]
\centering
\includegraphics[width=\textwidth]{image_lj_vs_threat.png}
\caption{Dual-axis visualization of the adaptive ratcheting window $L_j$ (orange solid, left axis) and composite threat level $t$ (red dashed, right axis) over 50 federated rounds. At round~1 (threat $t=0.000$), $L_j$ rises from the default of 10 to $L_{max}=20$ during the benign phase. At round~3, the peak threat of $t=0.879$ triggers an immediate contraction to $L_j=8$. As exponential decay dissipates the threat evidence, $L_j$ gradually recovers to $17$ (round~11), $18$ (round~19), and $19$ (round~20), stabilizing near $L_{max}$ for rounds $20$--$50$. The conservative asymmetry---fast contraction, slow expansion---directly implements the safety property of the dual-threshold ratchet policy.}
\label{fig:L_j_threat_combined}
\end{figure*}

The complete $L_j$ adjustment audit log is presented in Table~\ref{tab:L_j_audit}. The table records every transition of the adaptive ratcheting window alongside the concurrent composite threat score computed by the \texttt{ThreatMonitor}. The protocol executed a total of \textbf{6 asymmetric ratchet boundary adjustments} over 50 rounds, demonstrating the sparse, event-driven nature of the adaptive mechanism.

\begin{table*}
\centering
\caption{$L_j$ Adaptive Adjustment Audit Log (50-Round Prototype)}
\label{tab:L_j_audit}
\begin{tabular}{c c c c l}
\toprule
\textbf{Round} & $L_j^{prev}$ & $L_j^{new}$ & \textbf{Threat} $t$ & \textbf{Protocol Response} \\
\midrule
1  & 10 & 20 & 0.0000 & Benign start; rise to $L_{max}$ \\
3  & 20 & 8  & 0.8789 & Peak threat; immediate contraction \\
4  & 8  & 9  & 0.8478 & Threat persists; marginal relaxation \\
11 & 9  & 17 & 0.5359 & Significant decay; window expansion \\
19 & 17 & 18 & 0.4495 & Continued decay; incremental rise \\
20 & 18 & 19 & 0.4094 & Near-benign; approaching $L_{max}$ \\
\midrule
\multicolumn{2}{l}{\textbf{Rounds 20--50}} & $\approx$19 & $\approx$0.30 & Stable; residual decay plateau \\
\bottomrule
\end{tabular}
\end{table*}

\subsection{Zero-Threat Baseline Comparison}
\label{subsec:zero_threat}
To isolate the performance advantage of adaptive ratcheting under benign operational conditions, we conducted an additional controlled benchmark using 40 federated rounds with 4 clients and 200 samples per client. The \texttt{ThreatMonitor} received no security events throughout the run, yielding a composite threat score of $t = 0.0$ for every round. Under this zero-threat condition, the adaptive policy sets $L_j = L_{max} = 20$ (Eq.~\eqref{eq:adaptive_policy}), while the baseline PQBFL maintains a fixed $L_j = 10$. The results are summarised in Table~\ref{tab:zero_threat_comparison}.

\begin{table}[H]
\centering
\caption{Zero-Threat Benchmark: Baseline PQBFL vs.\ Adaptive PQBFL (40 rounds, 4 clients)}
\label{tab:zero_threat_comparison}
\begin{tabular}{l c c c}
\toprule
\textbf{Metric} & \textbf{Baseline} & \textbf{Adaptive} & $\Delta$ \\
\midrule
Effective $L_j$ (zero-threat) & 10 (fixed) & \textbf{20} ($L_{max}$) & $+10$ rounds \\
Asymmetric ratchets (40 rds.) & \textbf{3}$\times$ & \textbf{1}$\times$ & $-2$ ratchets \\
Avg.\ round latency & 1.93\,ms & \textbf{0.90\,ms} & $-53\%$ \\
Cumulative overhead (40 rds.) & $\sim$77\,ms & $\sim$36\,ms & $\mathbf{-41}$\,ms saved \\
FL accuracy (avg.) & 0.823 & 0.823 & no change \\
\bottomrule
\multicolumn{4}{l}{\small Each asymmetric ratchet $\approx 20$\,ms (Kyber KEM + 3-way handshake).}
\end{tabular}
\end{table}

At zero threat, the adaptive protocol doubles the ratchet window from 10 to 20, reducing the number of expensive ML-KEM re-keying events from 3 to 1 over 40 rounds. The average per-round latency drops by 53\% (from 1.93\,ms to 0.90\,ms) and cumulative cryptographic overhead is reduced by 41\,ms, with zero degradation in FL accuracy. This confirms Property~2 (Non-Linear Noise Suppression) of the \texttt{AdaptiveRatchetPolicy} (Eq.~\eqref{eq:adaptive_policy}): under benign conditions the quadratic mapping ($\gamma = 2.0$) keeps $L_j$ near $L_{max}$, maximising throughput without sacrificing security.

\begin{figure*}[!t]
\centering
\includegraphics[width=\textwidth]{graph1_lj_over_rounds.png}
\caption{Ratcheting window $L_j$ over 40 FL rounds under zero-threat conditions. Baseline PQBFL (blue, constant $L_j = 10$) triggers 3 asymmetric ratchets at rounds 11, 21, and 31 (red dashed lines). Adaptive PQBFL (green, $L_j = 20 = L_{max}$) triggers only 1 ratchet at round 21 (grey dashed). The green shaded zone highlights the additional 10 rounds between ratchet events saved by the adaptive policy.}
\label{fig:graph1_lj_rounds}
\end{figure*}

\begin{figure*}[!t]
\centering
\includegraphics[width=\textwidth]{graph3_timing_comparison.png}
\caption{Combined performance comparison for the zero-threat scenario (40 rounds, 4 clients). \textit{Top}: Per-round latency showing the impact of each asymmetric ratchet event. \textit{Bottom-left}: Cumulative overhead with 41\,ms saved by the adaptive policy. \textit{Bottom-right}: Key metrics bar chart summarising the 53.3\% latency speedup at identical accuracy. Adaptive PQBFL consistently outperforms the static baseline by eliminating unnecessary PQ key-encapsulation overhead.}
\label{fig:graph3_timing}
\end{figure*}

\subsection{Performance Analysis}
The experiment confirms that the proposed adaptive framework responds dynamically and correctly. At round~3, \texttt{AdaptiveRatchetPolicy} computed:
\begin{equation}
L_j = \left\lfloor L_{max} - (L_{max} - L_{min}) \cdot t^{\gamma} \right\rceil
= \left\lfloor 20 - 18 \times (0.879)^{0.5} \right\rceil = 8,
\end{equation}
contracting the ratcheting window from $20$ to $8$ in a single round and increasing the asymmetric key rotation frequency by a factor of $20/8 = 2.5\times$. As threat evidence aged past the exponential decay horizon $\lambda$, the composite score decayed monotonically: $t = 0.848 \to 0.536 \to 0.449 \to 0.409$, with a long-run residual of $\approx 0.30$. The zero-threat performance advantages of this mechanism were previously visualised in Fig.~\ref{fig:graph3_timing}. Table~\ref{tab:blockchain_perf} summarises the measured blockchain transaction performance across all 1,062 on-chain operations.
\begin{table*}
\centering
\caption{Blockchain Transaction Performance Across 1,062 Transactions}
\label{tab:blockchain_perf}
\begin{tabular}{l l l l l}
\toprule
\textbf{Tx Type} & \textbf{Count} & \textbf{Avg (ms)} & \textbf{Min (ms)} & \textbf{Max (ms)} \\
\midrule
\texttt{deploy\_contract}  & 1    & 34.99 & 34.99 & 34.99 \\
\texttt{register\_project} & 1    & 12.84 & 12.84 & 12.84 \\
\texttt{register\_client}  & 10   & 10.31 & 7.95  & 20.31 \\
\texttt{publish\_task}     & 50   & 7.80  & 6.73  & 24.28 \\
\texttt{update\_model}     & 500  & 7.93  & 6.97  & 25.64 \\
\texttt{feedback\_model}   & 500  & 8.18  & 7.14  & 31.21 \\
\midrule
\textbf{All transactions}  & \textbf{1062} & \textbf{8.00} & \textbf{6.73} & \textbf{36.48} \\
\bottomrule
\end{tabular}
\end{table*}

The global model accuracy improved from $60.50\%$ at round~0 to $\mathbf{81.75\%}$ at round~50 ($+21.25\%$ absolute improvement), converging above $80\%$ within just 2 rounds. This steep early convergence, achieved despite the added cryptographic overhead of Phase~2 hybrid key establishment, confirms that the per-round AES-256-GCM encryption and KDF$_S$ chain derivation do not impede the underlying FedAvg optimization. Steady-state per-round blockchain transactions (\texttt{publish\_task}, \texttt{update\_model}, \texttt{feedback\_model}) settled consistently below $10$\,ms, demonstrating that the blockchain audit layer adds negligible latency relative to local model training.

\subsection{Security Hardening Verification}
The security dashboard ran a complete suite of \textbf{10} automated verification checks. All 10 passed with zero warnings and zero failures. Table~\ref{tab:security_checks} details the verified modules and their implementations.

\begin{table*}
\centering
\caption{Security Hardening Verification Report (10/10 Checks Passed)}
\label{tab:security_checks}
\begin{tabular}{c l l}
\toprule
\textbf{\#} & \textbf{Hardening Measure} & \textbf{Implementation} \\
\midrule
1  & Constant-time comparison     & \texttt{hmac.compare\_digest} \\
2  & AES-256-GCM AEAD             & \texttt{AESGCM} via \textit{cryptography} \\
3  & Random 12-byte nonces        & \texttt{os.urandom(12)} \\
4  & GCM tamper detection         & Tag rejection on byte modification \\
5  & Constant-time Kyber KEM      & \texttt{pqcrypto/liboqs} C bindings \\
6  & Kyber roundtrip verification & keygen $\to$ encap $\to$ decap match \\
7  & Constant-time ECDH           & OpenSSL secp256k1 \\
8  & Ethereum sig verification    & \texttt{verify\_signer} + HMAC digest \\
9  & Safe model serialization     & \texttt{numpy.npz}, \texttt{allow\_pickle=False} \\
10 & Adaptive ratcheting bounds   & $L_j \in [L_{min}, L_{max}]$ verified \\
\bottomrule
\end{tabular}
\end{table*}

\subsection{Comparison Analysis}
Table~\ref{tab:comparison_table} provides a comprehensive comparison of the proposed Adaptive PQBFL against key baselines across five evaluation dimensions. The proposed framework uniquely combines post-quantum security, adaptive post-compromise recovery, and side-channel resistance in a single unified protocol, achieving the best-in-class result on all security-critical metrics while maintaining competitive round latency.

\begin{table*}
\centering
\caption{Comparative Performance of FL Frameworks}
\label{tab:comparison_table}
\begin{tabular}{l c c c c c}
\toprule
\textbf{Framework} & \textbf{Accuracy} & \textbf{Latency} & \textbf{PCS Window} & \textbf{PQ-Safe} & \textbf{SC-Safe} \\
\midrule
Standard FedAvg        & 60.5\%$^\dagger$ & 28.1\,ms & $\infty$ & No  & No  \\
ECDH-Only BC-FL        & 81.2\%           & 31.2\,ms & $\infty$ & No  & No  \\
Static Kyber ($L_j=5$) & 81.5\%           & 62.1\,ms & 5 rds    & Yes & No  \\
Static Kyber ($L_j=2$) & 81.6\%           & 98.4\,ms & 2 rds    & Yes & No  \\
\textbf{Adaptive PQBFL (ours)} & \textbf{81.75\%} & \textbf{43.8\,ms} & \textbf{$L_j(t)\in[2,20]$} & \textbf{Yes} & \textbf{Yes} \\
\bottomrule
\multicolumn{6}{l}{$^\dagger$ Unprotected FedAvg degrades further under poisoning attacks.}
\end{tabular}
\end{table*}

The Adaptive PQBFL achieves $81.75\%$ final accuracy, marginally outperforming Static Kyber variants due to the reduced communication overhead when $L_j$ is large during benign rounds. The $29.4\%$ latency improvement over Static Kyber ($L_j=5$) is directly attributable to fewer asymmetric ratchet invocations: the adaptive protocol triggered only 6 $L_j$-boundary adjustments across 50 rounds (compared to $\lceil 50/5 \rceil = 10$ for static), saving approximately $4 \times 4.5 = 18$\,ms of Kyber KEM overhead. Critically, the adaptive protocol is the only evaluated framework to provide post-compromise security (\textbf{PCS}) with a variable, threat-responsive window \emph{and} pass all 10 side-channel hardening checks.

\subsection{Computational Analysis}
The end-to-end computational cost confirms the protocol's practical viability. Steady-state blockchain transactions complete at an average of $8.00$\,ms, with per-round symmetric KDF$_S$ chain advancement completing in $<1$\,ms (two fixed-length HMAC evaluations, $\mathcal{O}(1)$). Asymmetric ratchet events (Kyber-512 encapsulation + ECDH) incur approximately $4.5$--$5$\,ms per occurrence, amortized over $L_j$ rounds. Over the 50-round experiment, the adaptive framework triggered only $\mathbf{6}$ $L_j$ adjustments, compared to $\lceil50/2\rceil=25$ ratchets required by a static $L_j=2$ maximum-security baseline---an \textbf{76\% reduction} in PQC key-exchange overhead.

Space complexity is $\mathcal{O}(1)$: only 64 bytes of active key material ($RK_j$ and $CK_{i,j}$, each 32 bytes) are retained at any time, independent of participant count, session length, or model size. The \texttt{ThreatMonitor} evaluation at each round takes $\mathcal{O}(|\mathcal{E}|) \approx 10\,\mu s$, and the \texttt{AdaptiveRatchetPolicy} mapping takes $\approx 1\,\mu s$---both negligible relative to the $\sim 22$\,ms local training time. These results confirm that the Threat-Adaptive PQBFL framework is computationally efficient, spatially lightweight, and operationally superior to all evaluated baselines.

\section{Conclusion}
This paper presented the Threat-Adaptive Ratcheting Post-Quantum Blockchain Federated Learning (PQBFL) framework, the first unified protocol combining classical-PQ hybrid key establishment, blockchain-anchored gradient integrity, threat-adaptive symmetric/asymmetric ratcheting, and side-channel hardening into a single cohesive system for privacy-preserving healthcare analytics.

The prototype evaluation over 50 federated rounds with 10 clients and the \texttt{ml\_kem\_512} KEM backend demonstrated the core contributions empirically. The global model accuracy improved from $60.50\%$ to $\mathbf{81.75\%}$ ($+21.25\%$), confirming that the post-quantum cryptographic layer does not introduce convergence degradation. The adaptive ratcheting mechanism correctly contracted the session window from $L_j=20$ to $L_j=8$ within a single round upon detecting a peak threat score of $t=0.879$, and gradually recovered to $L_j \approx 19$ as the exponential decay dissipated the threat evidence---requiring only 6 asymmetric ratchet boundary adjustments over 50 rounds, a \textbf{76\% reduction} over a static maximum-security ($L_j=2$) baseline. All 1,062 blockchain transactions completed with an average latency of $\mathbf{8.00}$\,ms, confirming negligible on-chain overhead. The full security verification suite confirmed all 10 hardening measures active: constant-time HMAC comparison, randomised AES-GCM nonces, GCM tamper detection, dynamic KDF salts, and Kyber/ECDH constant-time bindings.

The formal security analysis established three key properties of the \texttt{AdaptiveRatchetPolicy}: (i)~monotonicity ($t_1 > t_2 \Rightarrow L_j(t_1) \leq L_j(t_2)$), (ii)~non-linear noise suppression ($\gamma < 1$ suppresses transient spikes), and (iii)~boundary safety ($L_j \in [L_{min}, L_{max}]$ always). The dual-threshold ratchet trigger ensures the conservative effective window $L_j^{eff} = \min(L_j^{adaptive}, L_j(t))$ is enforced even during cooldown periods, eliminating race conditions between threat scoring and key rotation.

Future directions include the possibility of multi-authority blockchain governance of the decentralized consensus of ThreatMonitor, the incorporation of CRYSTALS-Dilithium (ML-DSA) of integration post-quantum gradient signing, and the future federation of implementation across geographically distributed nodes of the hospital networks with heterogeneous hardware goals. The framework provides a new security design floor in post-quantum, blockchain-based, adaptive federated learning in sensitive institutional settings.

\begin{thebibliography}{00}
\bibitem{kalapaaking2023} A. P. Kalapaaking, I. Khalil, and X. Yi, ``Blockchain-based Federated Learning with SMPC Model Verification Against Poisoning Attack,'' \textit{arXiv:2304.13360}, 2023.
\bibitem{sameera2024} K. M. Sameera et al., ``PRIVACY-PRESERVING IN BLOCKCHAIN-BASED FEDERATED LEARNING SYSTEMS,'' \textit{arXiv:2401.03552}, 2024.
\bibitem{yang2024_survey1} Z. Yang et al., ``A Survey and Comparison of Post-quantum and Quantum Blockchains,'' \textit{IEEE Communications Surveys}, 2024.
\bibitem{li2024} P. Li, T. Chen, and J. Liu, ``Enhancing Quantum Security over Federated Learning via Post-Quantum Cryptography,'' \textit{arXiv:2409.04637}, 2024.
\bibitem{baseri2025} Y. Baseri et al., ``Blockchain Security Risk Assessment in Quantum Era, Migration Strategies and Proactive Defense,'' \textit{arXiv:2501.11798}, 2025.
\bibitem{gharavi2025} H. Gharavi, J. Granjal, and E. Monteiro, ``PQBFL: A Post-Quantum Blockchain-based Protocol for Federated Learning,'' \textit{arXiv:2502.14464}, 2025.
\bibitem{commey2025_pqsbfl} D. Commey and G. V. Crosby, ``PQS-BFL: A Post-Quantum Secure Blockchain-based FL Framework,'' \textit{arXiv:2505.01866}, 2025.
\bibitem{yang2023_secureagg} S. Yang et al., ``A Post-quantum Secure Aggregation for Federated Learning,'' \textit{ACM/ICCNS}, 2023.
\bibitem{anon_survey_qsafe} Anonymous, ``A survey on quantum-safe blockchain security infrastructure,'' \textit{Journal of Information Security}, 2025.
\bibitem{choudhury2025} B. Choudhury et al., ``A Comprehensive Survey on Pre Versus Post Quantum Security Schemes for 5G-Enabled IoT,'' \textit{IEEE Access}, 2025.
\bibitem{kappala2026} A. Kappala et al., ``A Dynamic Quantum-Resistant Selective Encryption Approach for Agricultural Sensors,'' \textit{IEEE Access}, 2026.
\bibitem{nguyen2025} V.-D. Nguyen et al., ``A Novel Blockchain-Enabled Federated Learning Scheme for IoT Anomaly Detection,'' \textit{IEEE TMLCN}, 2025.
\bibitem{yang2024_survey2} Z. Yang et al., ``A Survey and Comparison of Post-Quantum and Quantum Blockchains (Vol 2),'' \textit{IEEE}, 2024.
\bibitem{marzougui2025} F. Marzougui et al., ``An Overview of Blockchain-Enabled Federated Learning Architectures for Agricultural Applications,'' \textit{Procedia Computer Science}, 2025.
\bibitem{bentley2014} R. A. Bentley et al., ``Big data in the new media environment,'' \textit{Journal of Archaeology}, 2014.
\bibitem{nezhadsistani2025} N. Nezhadsistani et al., ``Blockchain-Enabled Federated Learning in Healthcare,'' \textit{IEEE Access}, 2025.
\bibitem{dincer2022} H. Dincer, ``CRYPTOGRAPHIC PROTOCOLS OF SIGNAL AND SIGNAL BASED INSTANT MESSAGING APPLICATIONS,'' \textit{METU Thesis}, 2022.
\bibitem{zhang2025} Y. Zhang et al., ``Efficient Full-Stack Private Federated Deep Learning With Post-Quantum Security,'' \textit{IEEE TDSC}, 2025.
\bibitem{lopez2025} J. Lopez et al., ``Evaluating Post-Quantum Cryptographic Algorithms on Resource-Constrained Devices,'' \textit{arXiv:2507.08312}, 2025.
\bibitem{rong2025} C. Rong et al., ``Federated Large Domain Model System,'' \textit{Blockchain: Research and Applications}, 2025.
\bibitem{elbizri2026} M. El Bizri et al., ``Institutional Approaches to Post-Quantum Cryptography: Migration Frameworks,'' \textit{IEEE Access}, 2026.
\bibitem{chaudhary2024} D. Chaudhary et al., ``Module Lattice-Based Post Quantum Secure Blockchain Authentication Framework,'' \textit{IEEE Access}, 2024.
\bibitem{aygul2025} M. A. Aygul et al., ``Multi-User Secret Key Generation for WSNs via Wireless Channels,'' \textit{IEEE OJCOMS}, 2025.
\bibitem{chen2021} J. Chen et al., ``On the construction of a post-quantum blockchain for smart city,'' \textit{JISA}, 2021.
\bibitem{commey2025_health} D. Commey et al., ``Post-Quantum Secure Blockchain-Based Federated Learning Framework for Healthcare Analytics,'' \textit{IEEE}, 2025.
\bibitem{anon_pdf1} Anonymous, ``Advances in Post-Quantum Signatures,'' \textit{Independent Research}, 2025.
\bibitem{jose2021} J. M. Jose and P. V, ``A Survey on Consensus Algorithms in Blockchain based on Post Quantum Cryptosystems,'' \textit{IIIT Kottayam}, 2021.
\bibitem{gurung2024} D. Gurung et al., ``Performance analysis and evaluation of postquantum secure blockchained FL,'' \textit{Computer Networks}, 2024.
\bibitem{sharma2022} P. K. Sharma et al., ``Blockchain and Federated Learning-enabled Distributed Secure Architecture,'' \textit{ASSS}, 2022.
\bibitem{anon_crdownload} Anonymous, ``Investigating Smart Contract Vulnerabilities in Quantum Blockchain Environments,'' \textit{Tech Report}, 2024.
\bibitem{moore2024} E. Moore et al., ``A Survey on Secure and Private Federated Learning Using Blockchain,'' \textit{IEEE}, 2024.
\bibitem{eren2025} H. Eren et al., ``Security and Privacy in the Internet of Everything (IoE): A Review,'' \textit{Applied Sciences}, 2025.
\bibitem{javed2025} M. S. Javed et al., ``AI-Driven Blockchain and Federated Learning for Secure EHR Sharing,'' \textit{Electronics}, 2025.
\bibitem{wang2026} X. Wang et al., ``Securing Federated Learning With Blockchain in the Medical Field,'' \textit{JMIR}, 2026.
\bibitem{saeed2024} M. M. Saeed and F. Alqahtani, ``AI and post-quantum cryptography powered cybersecurity approaches for IoT systems,'' \textit{PeerJ CS}, 2024.
\bibitem{hlatshwayo2026} M. Q. Hlatshwayo et al., ``A Technical Review of Quantum Computing Use Cases for Finance and Economics,'' \textit{Quantum Reports}, 2026.
\end{thebibliography}

\end{document}
